\documentclass[conference]{IEEEtran}
\IEEEoverridecommandlockouts
\usepackage[utf8]{inputenc}
\usepackage[T1]{fontenc}
\usepackage[spanish,es-tabla,es-noquoting,es-noshorthands]{babel}
\usepackage{cite}
\usepackage{amsmath,amssymb,amsfonts}
\usepackage{booktabs,tabularx,multirow,array}
\usepackage{graphicx}
\usepackage{xcolor}
\usepackage{tikz}
\usetikzlibrary{arrows.meta,positioning,shapes.geometric,calc,fit,backgrounds,decorations.pathreplacing,babel}
\usepackage[colorlinks=true,urlcolor=blue!55!black,citecolor=blue!55!black,linkcolor=blue!55!black]{hyperref}

\newcolumntype{Y}{>{\raggedright\arraybackslash}X}
\newcolumntype{C}{>{\centering\arraybackslash}X}
\newcommand{\pct}{\mbox{\symbol{37}}}

\renewcommand{\topfraction}{0.95}
\renewcommand{\bottomfraction}{0.95}
\renewcommand{\textfraction}{0.05}
\renewcommand{\floatpagefraction}{0.88}
\renewcommand{\dbltopfraction}{0.95}
\renewcommand{\dblfloatpagefraction}{0.88}
\setcounter{topnumber}{4}
\setcounter{bottomnumber}{4}
\setcounter{totalnumber}{8}
\setcounter{dbltopnumber}{4}
\setlength{\textfloatsep}{4pt plus 1pt minus 1pt}
\setlength{\floatsep}{3pt plus 1pt minus 1pt}
\setlength{\dbltextfloatsep}{4pt plus 1pt minus 1pt}
\setlength{\dblfloatsep}{3pt plus 1pt minus 1pt}
\setlength{\abovecaptionskip}{2pt}
\setlength{\belowcaptionskip}{0pt}
\AtBeginDocument{%
  \setlength{\abovedisplayskip}{3.5pt plus 1pt minus 1pt}%
  \setlength{\belowdisplayskip}{3.5pt plus 1pt minus 1pt}%
  \setlength{\abovedisplayshortskip}{2pt plus 1pt minus 1pt}%
  \setlength{\belowdisplayshortskip}{2pt plus 1pt minus 1pt}%
}

\begin{document}

\title{EvidenCite: Recomendación Local de Citas y Clasificación de su Función Retórica en Artículos Académicos mediante Recuperación Densa e Inferencia en el Símplex}

\author{
\IEEEauthorblockN{Daniel R. Gallo Acevedo, Dubán L. Jiménez Puerta, Julián A. Malaver Montoya y Javier M. Torres Ortega}
\IEEEauthorblockA{\textit{Maestría en Inteligencia Artificial (MAIA), Departamento de Ingeniería de Sistemas y Computación}\\
\textit{Universidad de los Andes}, Bogotá D.C., Colombia}
}

\maketitle

\begin{abstract}
La clasificación automática de la función retórica de citas científicas suele operar solo sobre una oración del artículo citante o representar al documento citado mediante su título y resumen global, omitiendo la evidencia específica del cuerpo citado y reduciendo la taxonomía a esquemas gruesos de 3 a 6 etiquetas. Presentamos \textbf{EvidenCite}, un marco metodológico y experimental que integra la recuperación densa intra-documento de pasajes locales con la clasificación probabilística de intenciones citacionales sobre las 9 clases canónicas de Lyu et al. en el símplex $\Delta^8$. Mediante una arquitectura desacoplada en cinco capas, construimos un corpus de $18.658$ instancias con paridad estricta 1:1 entre la representación global (\textbf{Configuración A}: Título + Resumen) y la evidencia local recuperada (\textbf{Configuración B}: Top-3 fragmentos del cuerpo $\le 110$ palabras codificados con \textit{SciBERT} y sus encabezados de sección originales), con fuga cero entre particiones agrupadas por citante ($\mathcal{D}_{\text{train}}=13.077$, $\mathcal{D}_{\text{dev}}=2.783$, $\mathcal{D}_{\text{test}}=2.798$). Formalizamos geométricamente la triangulación retórica de la tríada metodológica (\texttt{Basis}, \texttt{Improvement}, \texttt{Comparison}) en ventanas de $m=3$ oraciones para calibrar el filtrado de supervisión débil NLI ($\tau=0,6, \theta_{\text{conf}}=0,50, \theta_{\text{margin}}=0,15$) y la síntesis determinista de demostraciones \textit{Few-Shot}. Evaluando un diseño factorial $2 \times 2 \times 2$ ($22.384$ inferencias de LLMs de frontera) sobre $\mathcal{D}_{\text{test}}$ ($N=2.798$ instancias curadas manualmente al $100\pct{}$ por humanos), demostramos que \texttt{Gemini-3.8-Flash} en régimen \textit{0-Shot} con \textbf{Configuración B} lidera globalmente y domina la frontera de Pareto ($\text{Macro-F1}=0,5996$, $\text{Weighted-F1}=0,6933$, $\text{Accuracy}=68,41\pct{}$, $\text{ECE}=0,1158$, $\$17,68\text{ USD}/10\text{k}$ citas), mientras la Configuración B aporta mejoras significativas en \texttt{GPT-4.1-mini} ($+2,86\text{ pp}$ global, McNemar $p=6,75 \times 10^{-6}$, y $+10,28\text{ pp}$ en \textit{Related Work}). Finalmente, la curaduría humana refuta el supuesto techo epistémico del panel 7B--8B, logrando $69,11\pct{}$ de exactitud en el estrato de desacuerdo previo y rescatando el $57,66\pct{}$ de los fallos del panel.
\end{abstract}

\begin{IEEEkeywords}
Clasificación de función de citas, recuperación densa de pasajes, encuadre retórico, modelos de lenguaje de gran escala, supervisión débil, calibración probabilística, cienciometría computacional.
\end{IEEEkeywords}

\section{Introducción y Planteamiento del Problema}

Las citas bibliográficas constituyen la estructura retórica mediante la cual un estudio científico se articula con el conocimiento previo~\cite{moravcsik1975some,lyu2021classification}. No obstante, los conteos bibliométricos tradicionales tratan toda referencia como un voto homogéneo, sin distinguir si provee contexto general (\texttt{Background}), identifica una brecha (\texttt{Gap}), aporta el fundamento teórico directo (\texttt{Basis}), provee una herramienta o dataset ejecutado directamente (\texttt{Application}), inspira una extensión metodológica (\texttt{Improvement}), respalda un hallazgo empírico (\texttt{Evidence}), acredita la autoría original (\texttt{Originator}), sugiere lecturas adicionales (\texttt{Further Reading}) o actúa como línea base de contraste (\texttt{Comparison})~\cite{lyu2021classification,zhang2025deep}. La Clasificación de Función de Citas (\textit{Citation Function Classification}, CFC) automatiza la identificación del rol comunicativo de cada mención en el discurso académico~\cite{teufel2006automatic,jurgens2018measuring}.

Pese al avance de los codificadores científicos preentrenados~\cite{beltagy2019scibert,cohan2019structural,paolini2025citefusion}, persisten tres cuellos de botella estructurales. Primero, mientras las síntesis bibliométricas identifican nueve funciones científicas observables en texto~\cite{lyu2021classification}, los benchmarks de PLN suelen colapsar la tarea a 3--6 categorías (SciCite~\cite{cohan2019structural}, ACL-ARC~\cite{jurgens2018measuring}, ACT~\cite{pride2020authoritative}) y recortar el contexto a una sola oración, perdiendo relaciones argumentativas multi-oración~\cite{lauscher2022multicite}. Segundo, los clasificadores existentes operan solo sobre el fragmento citante ($\mathcal{P}_{\text{citing}}$) o —como en recomendación local de citas en dos etapas~\cite{gu2022local} y codificadores por tripletas~\cite{singh2023scirepeval}— representan al artículo citado ($\mathcal{P}_{\text{cited}}$) mediante su resumen global (\textbf{Configuración A}: Título + Resumen), el cual es invariante entre menciones y omite los pasajes específicos del cuerpo que desambiguan funciones metodológicas cercanas. Tercero, aunque los Modelos de Lenguaje de Gran Escala (LLMs) mitigan la escasez de etiquetas~\cite{ding2023gpt3}, los modelos abiertos ligeros (7B--8B) y los LLMs sin calibración sufren de sobre-atribución metodológica (\textit{classifier sycophancy}), errores de atribución de terceros y baja concordancia con expertos humanos~\cite{fogelson2025llms}.

Para superar estas limitaciones, proponemos \textbf{EvidenCite}, un marco extremo a extremo que integra recuperación densa intra-documento de evidencia local en $\mathcal{P}_{\text{cited}}$ con clasificación probabilística sobre las 9 clases canónicas de Lyu et al.~\cite{lyu2021classification} en el símplex $\Delta^8$. Frente a la representación global (\textbf{Configuración A}), nuestra propuesta (\textbf{Configuración B}) segmenta el cuerpo de $\mathcal{P}_{\text{cited}}$ —extraído de \textit{unarXive}~\cite{besrour2026unarxive} y \textit{S2ORC}~\cite{lo2020s2orc}— en fragmentos oracionales ($\le 110$ palabras) y recupera con \textit{SciBERT}~\cite{beltagy2019scibert,reimers2019sentence} los $K=3$ pasajes más afines a la ventana citante ($m=3$ oraciones), preservando sus encabezados de sección autorales como dispositivos de encuadre (\textit{framing devices})~\cite{paolini2025citefusion,goffman1974frame,bertin2014study,pimentel2021disambiguatory}. Formulamos tres preguntas de investigación:
\begin{itemize}
    \item \textbf{PI1 (Evidencia Local y Encuadre):} ¿En qué medida y bajo qué ubicaciones retóricas los Top-3 fragmentos locales del citado con sus encabezados (\textbf{Config B}) superan al resumen global (\textbf{Config A}) bajo paridad 1:1?
    \item \textbf{PI2 (Triangulación en $\Delta^8$ y Supervisión Multi-Etapa):} ¿Cómo impacta la co-localización de la tríada metodológica (\texttt{Basis}, \texttt{Improvement}, \texttt{Comparison}) en ventanas de $m=3$ oraciones al filtrado NLI (\textit{DeBERTa-v3})~\cite{he2023debertav3}, al panel de jueces abiertos en \textit{vLLM}~\cite{kwon2023efficient} y a la validación humana?
    \item \textbf{PI3 (Frontera de Pareto, Calibración y Despliegue):} ¿Cuál es el balance entre Macro-F1, calibración (\textit{ECE}, \textit{Brier}), latencia, razonamiento latente y costo neto (bajo \textit{Prompt Caching}) al contrastar LLMs comerciales (\textit{0-Shot} vs. \textit{9-Shot}) frente a codificadores locales~\cite{beltagy2019scibert,hu2022lora}?
\end{itemize}

Nuestras contribuciones principales son: (1) un \textbf{corpus pareado 1:1 de $18.658$ instancias} en cinco capas con fuga cero ($\text{leakage}=0$), homologado a las 9 clases de Lyu et al.~\cite{lyu2021classification,peroni2012fabio} y evaluado sobre un benchmark de $N=2.798$ citas revisado manualmente al $100\pct{}$ por humanos; (2) el \textbf{teorema de co-localización y vecindad retórica en $\Delta^8$}, que fundamenta la calibración NLI ($\tau=0,6, \theta_{\text{conf}}=0,50, \theta_{\text{margin}}=0,15$) y la síntesis determinista \textit{9-Shot} (reduciendo en $4\times$ los ceros exactos); (3) una \textbf{evaluación factorial de $22.384$ inferencias} donde \texttt{Gemini-3.8-Flash · 0-Shot · Config B} alcanza el máximo global ($\text{Macro-F1} = 0,5996$, $\text{Weighted-F1} = 0,6933$, $\text{Accuracy} = 68,41\pct{}$, $\text{ECE} = 0,1158$, $\$17,68/10\text{k}$), la Config B mejora significativamente a \texttt{GPT-4.1-mini} ($+2,86\text{ pp}$ global, $p = 6,75 \times 10^{-6}$, y $+10,28\text{ pp}$ en \textit{Related Work}), y la curaduría humana refuta el techo epistémico en \textit{Needs Review} ($69,11\pct{}$ de exactitud en Gemini, rescatando el $57,66\pct{}$ de los fallos del panel 7B--8B); y (4) la \textbf{caracterización operativa} (supresión del $99,5\pct{}$ de tokens de pensamiento en 9-Shot y ahorro del $59,0\pct{}\text{--}60,8\pct{}$ por caché en OpenAI) junto con un demostrador cloud en tiempo real sobre AWS ECS Fargate.

\section{Estado del Arte y Revisión Bibliográfica}

\subsection{Evolución de Esquemas Taxonómicos y Complejidad Discursiva}

Desde la dicotomía conceptual-operacional de Moravcsik y Murugesan~\cite{moravcsik1975some} y el esquema computacional de 12 clases de Teufel et al.~\cite{teufel2006automatic}, la literatura se bifurcó entre taxonomías bibliométricas ricas como CiTO~\cite{peroni2012fabio} y benchmarks de PLN simplificados para elevar el acuerdo inter-anotador~\cite{zhang2025deep}: ACL-ARC ($1.941$ instancias, 6 clases)~\cite{jurgens2018measuring}, SciCite ($11.020$ instancias, 3 clases)~\cite{cohan2019structural} y ACT ($11.233$ citas, 6 clases)~\cite{pride2020authoritative}. Dos estudios expusieron sus límites: Lauscher et al.~\cite{lauscher2022multicite} introdujeron \textit{MultiCite} ($12.653$ contextos en $1.193$ artículos de S2ORC~\cite{lo2020s2orc}, 7 clases, acuerdo humano $0,70$), demostrando que las citas abarcan múltiples oraciones con alta co-ocurrencia pragmática (máxima información mutua entre \texttt{Extends} y \texttt{Differences}) y que una ventana de $m=3$ oraciones optimiza el F1 de \textit{SciBERT} frente a $m=1$ o $m \ge 5$. Paralelamente, la meta-síntesis de Lyu et al.~\cite{lyu2021classification} sobre 38 estudios (1965--2018) destiló las \textbf{9 clases científicas observables en texto} adoptadas en EvidenCite (\texttt{Background}, \texttt{Gap}, \texttt{Basis}, \texttt{Comparison}, \texttt{Application}, \texttt{Improvement}, \texttt{Evidence}, \texttt{Originator}, \texttt{Further Reading}), reconciliando granularidad cienciométrica y modelamiento computacional~\cite{zhang2025deep}.

\subsection{Recuperación Local de Citas y Representación Densa de Documentos Científicos}

\textit{SciBERT}~\cite{beltagy2019scibert} ($1,14\text{M}$ de artículos, $3,17\times 10^9$ tokens, vocabulario \textit{SciVocab} con $42\pct{}$ de solapamiento con BERT base) alcanzó $70,98$ Macro-F1 en ACL-ARC y $85,49$ en SciCite, mientras Reimers y Gurevych~\cite{reimers2019sentence} demostraron la superioridad del \textit{mean pooling} enmascarado para similitud densa sin ajuste fino. En representación documental, Singh et al.~\cite{singh2023scirepeval} presentaron \textit{SciRepEval} y \textit{SPECTER2} ($6,2\text{M}$ de tripletas en 23 disciplinas), advirtiendo que comprimir un artículo en un único vector estático de \texttt{Título + Resumen} carece de resolución para tareas intra-documento. En recomendación local de citas, Gu et al.~\cite{gu2022local} propusieron un pipeline de \textit{prefetching} (\textit{HAtten}) y re-ranking cruzado con \textit{SciBERT} ($\texttt{[CLS]} \; \text{Contexto Citante} \; \texttt{[SEP]} \; \text{Título} \parallel \text{Resumen Citado} \; \texttt{[SEP]}$), notando la elipsis de contextos de una sola oración. Aunque \textit{unarXive 2024}~\cite{besrour2026unarxive} escala a $2,28\text{M}$ de artículos completos extraídos de \LaTeX{} con $39,5\text{M}$ de citas resueltas y secciones preservadas, carece de etiquetas retóricas y el re-ranker de Gu et al.~\cite{gu2022local} sigue limitando $\mathcal{P}_{\text{cited}}$ a su \texttt{Título + Resumen} (\textbf{Config A}). EvidenCite traslada la recuperación densa al cuerpo completo de $\mathcal{P}_{\text{cited}}$ (\textbf{Config B}) para extraer los Top-3 fragmentos locales que sustentan la cita.

\subsection{Encuadre Retórico, Supervisión Débil y Retos con Modelos de Lenguaje de Gran Escala}

Apoyándose en la teoría de marcos de Goffman~\cite{goffman1974frame}, la estructura léxica IMRaD de Bertin y Atanassova~\cite{bertin2014study} y la ganancia informacional de prefijos de Pimentel et al.~\cite{pimentel2021disambiguatory}, Paolini et al.~\cite{paolini2025citefusion} demostraron en \textit{CiteFusion} que anteponer el encabezado de sección (\textit{framing device}, \textit{WS} vs. \textit{WoS}) eleva el Macro-F1 en SciCite ($89,60\pct{}$ vs. $88,48\pct{}$) y ACL-ARC ($76,24\pct{}$ vs. $71,46\pct{}$). Su análisis SHAP reveló que \textit{SciBERT} asigna alta atribución a los tokens de sección, que \texttt{Background} se clasifica por exclusión y que encabezados citantes genéricos (\textit{Introduction}, \textit{Related Work}) introducen ambigüedad cuando albergan citas metodológicas. De modo afín, \textit{CitePrompt}~\cite{lahiri2023citeprompt} aplicó \textit{prompt-tuning} con verbalizadores por sección de S2ORC ($68,39\pct{}$ en ACL-ARC y $86,33\pct{}$ en SciCite supervisado), pero colapsó en régimen \textit{Zero/Few-Shot} de 6 clases confundiendo \texttt{Extends} y \texttt{Background}.

Para escalar la anotación, Ding et al.~\cite{ding2023gpt3} mostraron que los LLMs reducen en $10\times\text{--}26\times$ el costo de etiquetado aunque pierden precisión al crecer las clases; paralelamente, \textit{DeBERTa-v3}~\cite{he2023debertav3} habilita filtrado NLI calibrado, y \textit{vLLM}~\cite{kwon2023efficient} junto con \textit{LoRA}~\cite{hu2022lora} viabilizan comités y ajuste eficiente de modelos abiertos. Sin embargo, Fogelson et al.~\cite{fogelson2025llms} evaluaron modelos de 7B--8B (\texttt{Llama-3.1-8B}, \texttt{Mistral-7B}) y \texttt{GPT-4.1-mini} en ventanas de 3 oraciones ($1,27\text{M}$ de citas, 3 clases), reportando baja concordancia inter-LLM ($\text{F1} = 0,49$), rezago de modelos 7B--8B ($\text{F1} = 0,50$ vs. $0,67$ en GPT) y dos sesgos críticos: \textit{atribución errónea de terceros} en revisiones de literatura y \textit{adulación del clasificador} (\textit{classifier sycophancy}, sobre-predicción metodológica). La Tabla~\ref{tab:sota_comparativa} contrasta la literatura frente a \textbf{EvidenCite}.

\begin{table*}[!t]
\centering
\caption{Comparación Crítica del Estado del Arte en Clasificación y Recuperación Local de Citas frente a EvidenCite}
\label{tab:sota_comparativa}
\scriptsize
\setlength{\tabcolsep}{3.5pt}
\renewcommand{\arraystretch}{1.05}
\resizebox{\textwidth}{!}{%
\begin{tabular}{>{\raggedright\arraybackslash}p{2.7cm} >{\raggedright\arraybackslash}p{3.3cm} c c >{\raggedright\arraybackslash}p{3.2cm} >{\raggedright\arraybackslash}p{3.4cm} >{\raggedright\arraybackslash}p{5.4cm}}
\toprule
\textbf{Estudio / Sistema} & \textbf{Paradigma / Modelo Base} & \textbf{Corpus ($N$)} & \textbf{Clases} & \textbf{Ventana Citante y Encuadre} & \textbf{Evidencia del Citado ($\mathcal{P}_{\text{cited}}$)} & \textbf{Limitación Crítica Superada por EvidenCite} \\
\midrule
Jurgens et al.~\cite{jurgens2018measuring} (\textit{ACL-ARC}) & Características léxicas + Random Forest & $1.941$ (CL) & 6 & 1 oración; sin sección en entrada & Ninguna ($\varnothing$) & Contexto de 1 oración insuficiente; ausencia total de evidencia del artículo citado. \\
Cohan et al.~\cite{cohan2019structural} (\textit{SciCite}) & BiLSTM + ELMo + \textit{Structural Scaffolds} & $11.020$ (CS/Med) & 3 & 1 oración; sección como tarea auxiliar & Ninguna ($\varnothing$) & Taxonomía gruesa de 3 clases; pierde granularidad entre base, uso, mejora y evidencia. \\
Beltagy et al.~\cite{beltagy2019scibert} (\textit{SciBERT}) & Transformer 110M (\textit{SciVocab} completo) & ACL-ARC / SciCite & 6 / 3 & 1 oración aislada & Ninguna ($\varnothing$) & Ignora encabezados retóricos, contexto multi-oración y contenido del documento citado. \\
Lauscher et al.~\cite{lauscher2022multicite} (\textit{MultiCite}) & \textit{SciBERT} / \textit{RoBERTa} multi-etiqueta + QA & $12.653$ (CL) & 7 & Flexible / $m=3$ con \texttt{<cite>} & Ninguna en clasificación ($\varnothing$) & No incorpora pasajes del artículo citado; $93,6\pct{}$ de instancias sin sección estructurada. \\
Gu et al.~\cite{gu2022local} (\textit{HAtten+SciBERT}) & Prefetching \textit{HAtten} + Reranker \textit{SciBERT} & $1,66\text{M}$ arXiv / ACL & Ranking & $\pm 200$ car. + Título/Resumen citante & Global (\textbf{Config A}: Título + Resumen) & No clasifica intención retórica; representa al citado solo por su resumen global estático. \\
Singh et al.~\cite{singh2023scirepeval} (\textit{SPECTER2}) & \textit{SciBERT} + $6,2\text{M}$ tripletas + Adaptadores & SciRepEval (24 tareas) & 4 formatos & Nivel documental (sin contexto local) & Global (\textbf{Config A}: Título + Resumen) & Vector global invariante entre menciones; no aísla pasajes específicos del cuerpo citado. \\
Lahiri et al.~\cite{lahiri2023citeprompt} (\textit{CitePrompt}) & \textit{SciBERT} + \textit{OpenPrompt} + Verbalizador S2ORC & ACL-ARC / SciCite & 6 / 3 & 1 oración + plantilla prefijo MLM & Ninguna ($\varnothing$) & Colapsa en \textit{0/Few-Shot} sobre esquemas desbalanceados de 6 clases; confunde \texttt{Extends}/\texttt{Background}. \\
Paolini et al.~\cite{paolini2025citefusion} (\textit{CiteFusion}) & Parejas OVA (\textit{SciBERT}+\textit{XLNet}) + FFNN + SHAP & ACL-ARC / SciCite & 6 / 3 & 1 oración + prefijo sección citante (\textit{WS}) & Ninguna ($\varnothing$) & Encuadre solo en el citante (ambiguo en \textit{Intro}/\textit{Related Work}); sin evidencia del citado. \\
Fogelson et al.~\cite{fogelson2025llms} (\textit{LLM CIC}) & \texttt{Llama-3.1-8B}, \texttt{Mistral-7B}, \texttt{GPT-4.1-mini} & $1,27\text{M}$ (Eval $300$) & 3 & Ventana de $m=3$ oraciones & Ninguna ($\varnothing$) & Esquema grueso (3 clases); baja concordancia inter-LLM ($0,49$ F1) y \textit{sycophancy} sin calibración. \\
\midrule
\textbf{EvidenCite (Este Trabajo)} & \textbf{Recuperación \textit{SciBERT} Top-3 + Supervisión NLI/Panel/Gold + \textit{SciBERT}/\textit{LoRA} y LLMs API ($\Delta^8$)} & \textbf{$\mathbf{18.658}$ pares 1:1 ($\mathbf{2.798}$ Test Gold humano)} & \textbf{9 (Lyu et al.~\cite{lyu2021classification})} & \textbf{Ventana $\mathbf{m=3}$, \texttt{<cite>}, \texttt{[OTHERCIT]} + sección citante (\textit{Nearest Scope})} & \textbf{Paridad 1:1 Config A (Título+Resumen) vs. Config B (Top-3 Chunks $\le 110$ pal. + Sección Citada)} & \textbf{Integra taxonomía fina de 9 clases en $\Delta^8$, evidencia local recuperada del cuerpo citado, validación humana Gold A3, calibración y frontera de Pareto.} \\
\bottomrule
\end{tabular}%
}
\end{table*}

\section{Metodología}

\subsection{Visión General y Protocolo Experimental Paritario 1:1}

Formulamos la clasificación de función citacional como un problema de estimación probabilística condicionada sobre pares de documentos científicos $(\mathcal{P}_{\text{citing}}, \mathcal{P}_{\text{cited}})$. Sea $\mathcal{Y}_9 = \{c_1, \dots, c_9\}$ el espacio de las 9 funciones retóricas canónicas derivadas de la meta-síntesis de Lyu et al.~\cite{lyu2021classification}:
\begin{equation}
\label{eq:espacio_clases}
\begin{aligned}
\mathcal{Y}_9 = \{&\texttt{Background}, \texttt{Gap}, \texttt{Basis}, \texttt{Comparison},\\
&\texttt{Application}, \texttt{Improvement}, \texttt{Evidence},\\
&\texttt{Originator}, \texttt{Further Reading}\}.
\end{aligned}
\end{equation}
Toda arquitectura evaluada $f_\Theta: \mathcal{X}_\mu \to \Delta^8$ mapea una representación de entrada $\mathbf{x}_i^{(\mu)} \in \mathcal{X}_\mu$ bajo la modalidad de evidencia $\mu \in \{\text{Config A}, \text{Config B}\}$ hacia el símplex de probabilidad de dimensión ocho:
\begin{equation}
\label{eq:simplex_def}
\Delta^8 = \left\{ \mathbf{p} = (p_1, \dots, p_9)^\top \in [0, 1]^9 \;\middle|\; \sum_{k=1}^{9} p_k = 1 \right\},
\end{equation}
de donde la predicción puntual se obtiene como $\hat{y}_i = \arg\max_{c_k \in \mathcal{Y}_9} p_{i,k}$ con nivel de confianza asociado $\hat{\gamma}_i = \max_{k} p_{i,k}$.

Como ilustra la Figura~\ref{fig:conceptual_citante_citado}, ambas configuraciones comparten idéntica información del artículo citante —la sección retórica $h_{\text{citing}}$ y la ventana de $m=3$ oraciones $x_{\text{ctx}}$— y difieren exclusivamente en la evidencia provista sobre el artículo citado $\mathcal{P}_{\text{cited}}$: la \textbf{Configuración A} suministra el título $t_{\text{cited}}$ y el resumen global $\text{abs}_{\text{cited}}$, mientras que la \textbf{Configuración B} inyecta los $K=3$ fragmentos locales $\mathcal{C}_{\text{top3}} = \{(c_{\pi_r}, h_{\pi_r}, s_{\pi_r})\}_{r=1}^3$ recuperados densamente del cuerpo de $\mathcal{P}_{\text{cited}}$ junto con sus encabezados de sección originales $h_{\pi_r}$. Para aislar causalmente el aporte de la evidencia local recuperada sin sesgos de muestreo, imponemos el \textbf{Invariante de Paridad Estricta 1:1}:
\begin{equation}
\label{eq:paridad_1_a_1}
\begin{aligned}
\mathcal{D}_s(\text{Config A}) &\equiv \mathcal{D}_s(\text{Config B}), \quad \forall s \in \{\text{train}, \text{dev}, \text{test}\},\\
\text{abs}_{\text{cited}}(i) \neq \varnothing \;&\land\; |\mathcal{C}_{\text{top3}}(i)| \ge 1, \quad \forall i \in \{1, \dots, N\}.
\end{aligned}
\end{equation}

\begin{figure}[!t]
\centering
\resizebox{0.92\columnwidth}{!}{%
\begin{tikzpicture}[
    font=\sffamily\scriptsize,
    container/.style={draw=black!65, rounded corners=2.5pt, fill=gray!6, inner sep=3.5pt},
    subboxA/.style={draw=blue!60!black, rounded corners=2pt, fill=blue!6, inner sep=2.5pt, align=left},
    subboxB/.style={draw=teal!70!black, rounded corners=2pt, fill=teal!7, inner sep=2.5pt, align=left},
    arrow/.style={-{Latex[length=2.0mm, width=1.4mm]}, thick}
]
\node[container, text width=3.7cm, minimum height=3.5cm, anchor=north west] (citing) at (0,0) {};
\node[anchor=north west, font=\sffamily\bfseries\scriptsize, text=black!85] at ([xshift=3pt,yshift=-2pt]citing.north west) {Artículo Citante $\mathcal{P}_{\text{citing}}$};
\node[draw=black!40, fill=white, rounded corners=2pt, text width=3.4cm, inner sep=2.5pt, anchor=north west, font=\sffamily\tiny] (csec) at ([xshift=3pt,yshift=-13pt]citing.north west) {
    \textbf{Sección ($h_{\text{citing}}$):} \textit{2. Related Work}\\[1.5pt]
    \textbf{Ventana ($m=3$):} $s_{i-1}$: Prior estimators incur exponential time \texttt{[OTHERCIT]}. $s_i$: We build upon \textcolor{blue!75!black}{\textbf{\texttt{<cite>Hopkins et al., 2022</cite>}}} and exploit their relaxation. $s_{i+1}$: We precondition covariance \texttt{[OTHERCIT]}.
};
\node[container, text width=4.3cm, minimum height=3.5cm, anchor=north west] (cited) at (4.5,0) {};
\node[anchor=north west, font=\sffamily\bfseries\scriptsize, text=black!85] at ([xshift=3pt,yshift=-2pt]cited.north west) {Artículo Citado $\mathcal{P}_{\text{cited}}$ (Paridad 1:1)};
\node[subboxA, text width=3.95cm, anchor=north west, font=\sffamily\tiny] (cfga) at ([xshift=3pt,yshift=-13pt]cited.north west) {
    \textbf{\textcolor{blue!70!black}{Config A (Global):}} Título $t_{\text{cited}} +$ Resumen $\text{abs}_{\text{cited}}$ (\textit{Efficient mean estimation with pure DP...}).
};
\node[subboxB, text width=3.95cm, anchor=north west, font=\sffamily\tiny] (cfgb) at ([xshift=3pt,yshift=-44pt]cited.north west) {
    \textbf{\textcolor{teal!70!black}{Config B (Top-3 Chunks $\le 110$ pal. + Sección):}}\\
    $\bullet$ R1 ($0,894$): \texttt{[2.2.1 Private mean]} \textit{``Our algorithm draws on...''}\\
    $\bullet$ R2 ($0,889$): \texttt{[6 Fine Estimation]} \textit{``SDP relaxation yields...''}\\
    $\bullet$ R3 ($0,888$): \texttt{[6 Fine Estimation]} \textit{``Median of means...''}
};
\draw[arrow, blue!70!black] ([yshift=10pt]csec.east) -- node[above, font=\sffamily\bfseries\tiny] {Par A} ([yshift=1pt]cfga.west);
\draw[arrow, teal!75!black] ([yshift=-12pt]csec.east) -- node[below, font=\sffamily\bfseries\tiny] {$T_{\text{query}}$ Top-3} ([yshift=2pt]cfgb.west);
\end{tikzpicture}%
}
\caption{Instancia pareada 1:1 en EvidenCite (inspirada en \textit{CiteFusion}~\cite{paolini2025citefusion}): evidencia global (\textbf{Config A}) vs. evidencia local Top-3 con encabezados (\textbf{Config B}). Elaboración propia.}
\label{fig:conceptual_citante_citado}
\end{figure}

\subsection{Pipeline Desacoplado de Datos en Cinco Capas e Ingesta Dirigida (A1)}

Para garantizar reproducibilidad determinista, trazabilidad de linaje y cero acoplamiento entre la extracción documental pesada y la experimentación de modelos, diseñamos una arquitectura secuencial en cinco capas ($\text{Capa 1: Bruta} \to \text{Capa 2: Estructurada} \to \text{Capa 3: Procesada} \to \text{Capa 4: Conjunto Pareado} \to \text{Capa 5: Etiquetada}$), diagramada en la Figura~\ref{fig:arquitectura_5_capas}.

\begin{figure*}[!t]
\centering
\resizebox{0.91\textwidth}{!}{%
\begin{tikzpicture}[
    font=\sffamily\scriptsize,
    chevron/.style={draw=black!70, fill=blue!10, rounded corners=2pt, minimum height=5.2mm, text width=3.15cm, align=center, font=\sffamily\bfseries\scriptsize},
    layerbox/.style={draw=black!60, rounded corners=2pt, fill=gray!4, text width=3.15cm, minimum height=2.25cm, inner sep=2.5pt, align=left, font=\sffamily\tiny},
    arrow/.style={-{Latex[length=2.0mm, width=1.4mm]}, semithick, black!75}
]
\node[chevron, fill=gray!14] (c1) at (0,0) {Capa 1: Ingesta Bruta};
\node[chevron, fill=blue!10] (c2) at (3.6,0) {Capa 2: Estructurada};
\node[chevron, fill=cyan!12] (c3) at (7.2,0) {Capa 3: Procesada};
\node[chevron, fill=teal!12] (c4) at (10.8,0) {Capa 4: Dataset 1:1};
\node[chevron, fill=green!12] (c5) at (14.4,0) {Capa 5: Etiquetada};
\draw[arrow] (c1.east) -- (c2.west); \draw[arrow] (c2.east) -- (c3.west); \draw[arrow] (c3.east) -- (c4.west); \draw[arrow] (c4.east) -- (c5.west);
\node[layerbox, anchor=north] (b1) at (0,-0.32) {\textbf{Fuentes Inmutables:}\\$\bullet$ \textbf{MultiCite:} $11.251$ citas (7 clases).\\$\bullet$ \textbf{unarXive:} Volcados \LaTeX{} CS.\\$\bullet$ \textbf{Caché Local:} Metadatos $O(1)$ y preprints \texttt{ar5iv}/PDF.};
\node[layerbox, anchor=north] (b2) at (3.6,-0.32) {\textbf{Corpus Estructurado:}\\$\bullet$ $\mathbf{66.314}$ artículos; $30.722$ citados con cuerpo.\\$\bullet$ \textbf{Seminal:} Top-250 $\to 146$ validados $\to 135$ cuerpos.\\$\bullet$ \textbf{Ámbito Envolvente} $\tilde{h}(p_t)$.};
\node[layerbox, anchor=north] (b3) at (7.2,-0.32) {\textbf{Unidades Atómicas:}\\$\bullet$ \textbf{Contextos ($m=3$):} $179.114 \to \mathbf{144.826}$ pares ($\bar{w}=83,94$ pal.), \texttt{<cite>}, \texttt{[OTHERCIT]}.\\$\bullet$ $3.629$ Gold + $4.445$ léxicas + $136.752$ pool.\\$\bullet$ \textbf{Chunks:} $W_{\max} \le 110$ pal.};
\node[layerbox, anchor=north] (b4) at (10.8,-0.32) {\textbf{Ensamble Pareado 1:1:}\\$\bullet$ \textbf{NLI ($\Delta^8$):} \textit{DeBERTa-v3} ($\tau=0,6$): $63.802 / 135.035$ ($47,25\pct{}$) $\to 10.584$.\\$\bullet$ \textbf{Top-3 \textit{SciBERT}:} $\bar{s}_{\pi_1}=0,8943$.\\$\bullet$ \textbf{Split ($\text{leakage}=0$):} $\mathbf{18.658}$ ($13.077 / 2.783 / 2.798$).};
\node[layerbox, anchor=north] (b5) at (14.4,-0.32) {\textbf{Auditoría A2 y Gold A3:}\\$\bullet$ \textbf{Panel 3 Jueces (vLLM):} $\bar{\kappa}=0,3939$; $A=4.772, B=2.667, C=4.705, D=2.885, F=3.629$.\\$\bullet$ \textbf{Curaduría A3:} $100\pct{}$ de $\mathcal{D}_{\text{test}}$ ($\mathbf{N=2.798}$ Gold, $54,32\pct{}$ corr.; $5.864$ Gold tot.).};
\draw[arrow] (b1.east) -- (b2.west); \draw[arrow] (b2.east) -- (b3.west); \draw[arrow] (b3.east) -- (b4.west); \draw[arrow] (b4.east) -- (b5.west);
\end{tikzpicture}%
}
\caption{Arquitectura desacoplada en cinco capas del pipeline de datos y etiquetado multi-etapa de EvidenCite (Actividades A1, A2 y A3). Elaboración propia.}
\label{fig:arquitectura_5_capas}
\end{figure*}

\subsubsection{Ingesta Dirigida y Resolución Bibliográfica Seminal}
En la Capa 1, el sistema ingesta las $11.251$ instancias humanas de \textit{MultiCite}~\cite{lauscher2022multicite} y los volcados \LaTeX{} de \textit{unarXive}~\cite{besrour2026unarxive} en Ciencias de la Computación. Cruzar directamente \textit{MultiCite} contra el subconjunto base de \textit{unarXive} resolvía el cuerpo citado en apenas $0,4\pct{}$ ($43/11.251$) de las instancias, lo que habría descartado $>99\pct{}$ del Gold humano bajo la Ec.~\eqref{eq:paridad_1_a_1}. Al observar que los $250$ artículos más citados en \textit{MultiCite} concentraban $\approx 35\pct{}$ de las menciones, ejecutamos un \textbf{pre-calentamiento seminal} con caché relacional y resolución en dos niveles: (i) HTML semántico vía \texttt{ar5iv} y (ii) conversión estructural de PDF en memoria ($\le 20$ páginas, $\le 20\text{ MB}$). De los 250 marcadores se validaron $146$ preprints en arXiv y se estructuraron $135$ cuerpos completos, consolidando en la Capa 2 un total de $\mathbf{66.314}$ artículos completos ($36.684$ citados indexados y $30.722$ con fragmentos válidos) y rescatando $\mathbf{3.629}$ instancias Gold con paridad 1:1.

\subsubsection{Preservación de Encabezados vía Ámbito Envolvente Más Próximo}
En vez de colapsar los títulos de sección hacia categorías rígidas IMRaD —perdiendo señales como \textit{``3.2 Self-Attention Layer''}—, la Capa 2 preserva la cadena autoral exacta. Para párrafos con encabezado nulo por artefactos \LaTeX{}, propagamos recursivamente sobre $(p_1, \dots, p_P)$ la regla de \textbf{Ámbito Envolvente Más Próximo} (\textit{Nearest Enclosing Scope}), fundamentada en los \textit{framing devices} de \textit{CiteFusion}~\cite{paolini2025citefusion,goffman1974frame}:
\begin{equation}
\label{eq:nearest_enclosing_scope}
\tilde{h}(p_t) = \begin{cases}
\text{clean}(h(p_t)), & \text{si } h(p_t) \neq \varnothing \land |\text{clean}(h(p_t))| > 0, \\
\tilde{h}(p_{t-1}), & \text{si } h(p_t) = \varnothing \land t > 1, \\
\text{``Preamble''}, & \text{si } h(p_1) = \varnothing.
\end{cases}
\end{equation}

\subsubsection{Extracción de Ventana Citante (\texorpdfstring{$m=3$}{m=3}) y Enmascaramiento de Co-Citas}
En la Capa 3a, cada sección $\mathcal{S}_k = (s_0, \dots, s_{n-1})$ se segmenta en oraciones protegiendo abreviaturas científicas (\textit{et al.}, \textit{i.e.}, \textit{Fig.}, decimales) y fórmulas. Siguiendo la optimalidad empírica de $m=3$ en \textit{MultiCite}~\cite{lauscher2022multicite}, para cada oración citante $s_i$ extraemos una ventana de tres oraciones con desplazamiento en bordes de sección:
\begin{equation}
\label{eq:ventana_m3}
\mathcal{W}_3(s_i; \mathcal{S}_k) = \begin{cases}
(s_0, \dots, s_{n-1}), & \text{si } n < 3, \\
(s_0, s_1, s_2), & \text{si } i = 0, \\
(s_{n-3}, s_{n-2}, s_{n-1}), & \text{si } i = n - 1, \\
(s_{i-1}, s_i, s_{i+1}), & \text{si } 0 < i < n - 1.
\end{cases}
\end{equation}
En $x_{\text{ctx}} = \bigoplus_{s \in \mathcal{W}_3(s_i; \mathcal{S}_k)} s$, la cita objetivo $c_{\text{target}}$ se delimita como \texttt{<cite>}$c_{\text{target}}$\texttt{</cite>} y toda co-cita secundaria $c_{\text{sec}} \neq c_{\text{target}}$~\cite{lyu2021classification} se reemplaza por \texttt{[OTHERCIT]} para impedir fugas de atribución retórica entre referencias adyacentes. En el corpus final, los contextos promedian $\bar{w} = 83,94$ palabras ($\min = 8$, $\max = 542$).

\subsubsection{Contratos Formales de Datos Invariantes entre Capas}
Para construir deterministamente las entradas en tiempo de ejecución según la configuración ($\mu \in \{\text{Config A}, \text{Config B}, \text{Solo Contexto}\}$), cada capa cumple un contrato algebraico invariante:
\begin{itemize}
    \item \textbf{Capa 2 (Artículo $\mathcal{P}$):} $\mathcal{S}_k = (k, \tilde{h}_k, (\sigma_{k,j})_{j=1}^{N_k})$ agrupa oraciones y $\mathcal{B}_{\text{bib}}$ mapea citas a identificadores (arXiv, DOI, OpenAlex):
    \begin{equation}
    \label{eq:contrato_capa2}
    \mathcal{P} = \left(p_{\text{id}}, s_{\text{id}}, t_{\text{title}}, \text{abs}, d_{\text{disc}}, \mathcal{G}_{\text{cat}}, (\mathcal{S}_k)_{k=0}^{S-1}, \mathcal{B}_{\text{bib}}\right).
    \end{equation}
    \item \textbf{Capa 3 (Contexto $\mathcal{X}_{\text{cite}}$ y Fragmento $c_{k,j}$):} $\mathcal{X}_{\text{cite}}$ vincula $(id_{\text{ctx}}, p_{\text{citing}}, h_{\text{citing}}, x_{\text{ctx}})$ con $(p_{\text{cited}}, t_{\text{cited}}, \text{abs}_{\text{cited}}, b_{\text{chunks}})$, y $c_{k,j} = (id_{\text{chunk}}, p_{\text{cited}}, k, \tilde{h}_k, j, \text{text}, n_{\text{tok}}, n_{\text{word}})$.
    \item \textbf{Capas 4 y 5 (Instancia Consolidada $\mathcal{I}$):}
    \begin{equation}
    \label{eq:contrato_capa4_5}
    \mathcal{I} = \left(id_{\text{inst}}, y, \mathcal{A}_{\text{meta}}, \mathcal{P}_{\text{citing}}^{\text{info}}, x_{\text{ctx}}, \mathcal{P}_{\text{cited}}^{\text{info}}\right),
    \end{equation}
    donde $\mathcal{P}_{\text{citing}}^{\text{info}} = (p_{\text{citing}}, s_{\text{id}}, t_{\text{citing}}, h_{\text{citing}})$, $\mathcal{P}_{\text{cited}}^{\text{info}} = (p_{\text{cited}}, t_{\text{cited}}, \text{abs}_{\text{cited}}, \mathcal{C}_{\text{top3}})$ reúne Config A y B, y $\mathcal{A}_{\text{meta}} = (d_{\text{src}}, y_{\text{orig}}, \mu_{\text{map}}, \ell_{\text{src}}, b_{\text{gold}}, P_1, \Delta,$ $\alpha_{\text{panel}}, y_{\text{panel}}, \mathcal{V}_{\text{panel}}, \ell_{\text{prior}}, s_{\text{split}})$ preserva el linaje multi-etapa (Figura~\ref{fig:json_contrato}).
\end{itemize}

\begin{figure}[!t]
\centering
\begin{tikzpicture}
\node[draw=black!60, fill=gray!5, rounded corners=2pt, inner sep=3.5pt, text width=0.95\columnwidth, font=\ttfamily\tiny] {
\{"instance\_id": "inst\_04783", "label": "Basis",\\
\ "annotation\_metadata": \{"source\_dataset": "unarxive\_mined", "label\_source": "human\_review", "is\_gold": true, "panel\_agreement": 0.67, "panel\_label": "Basis", "prior\_label\_source": "llm\_panel", "split": "test"\},\\
\ "citing\_paper": \{"id": "2212.08018", "rhetorical\_section": "High-level overview"\},\\
\ "citation\_context": "...procedure builds on <cite>Hopkins et al., 2022</cite>...",\\
\ "cited\_paper": \{"id": "2111.12981", "title": "Efficient mean estimation with pure DP...", "abstract": "We give the first polynomial-time algorithm...",\\
\ \ "retrieved\_chunks": [\{"chunk\_id": "2111.12981\_s4\_c1", "section\_heading": "2.2.1 Private mean...", "text": "Our algorithm draws heavily on...", "similarity\_score": 0.8936, "rank": 1\}]\}\}
};
\end{tikzpicture}
\caption{Contrato JSON invariante de Capas 4 y 5 ($\mathcal{D}_{\text{test}}$), encapsulando evidencia global (Config A), fragmentos locales con encabezados (Config B) y linaje multi-etapa. Elaboración propia.}
\label{fig:json_contrato}
\end{figure}

\subsection{Segmentación Granular y Recuperación Densa Top-3 con SciBERT (A1)}

\subsubsection{Algoritmo de Segmentación Acotada por Presupuesto de Tokens}
Los codificadores BERT (\textit{SciBERT}~\cite{beltagy2019scibert}) operan con ventana fija $L_{\max} = 512$ sub-palabras (\textit{WordPiece}). Como el segmento citante ocupa $60\text{--}110$ tokens, restan $\approx 400\text{--}440$ tokens para $\mathcal{P}_{\text{cited}}$; bloques de 300 palabras ($\approx 400$ tokens) truncarían los fragmentos de rango 2 y 3. Para alojar los $K=3$ pasajes sin fracturar oraciones ni cruzar fronteras de sección, dividimos cada sección $\mathcal{S}_k = (s_{k,1}, \dots, s_{k,N_k})$ de $\mathcal{P}_{\text{cited}}$ en fragmentos contiguos $c_{k,j} = \bigoplus_{i=a_j}^{b_j} s_{k,i}$ mediante acumulación voraz con cota $W_{\max} = 110$ palabras:
\begin{equation}
\label{eq:chunking_voraz}
\sum_{i=a_j}^{b_j} w(s_{k,i}) \le W_{\max} = 110 < \sum_{i=a_j}^{b_j + 1} w(s_{k,i}), \quad \text{para } b_j > a_j,
\end{equation}
permitiendo $a_j = b_j$ solo si una oración individual excede $W_{\max}$. Cada fragmento en Capa 3b promedia $\approx 108$ palabras ($\approx 140$ tokens), garantizando que los $K=3$ pasajes ($\approx 420$ tokens) encajen dentro de los 512 tokens.

\subsubsection{Codificación Densa por Mean Pooling y Consulta Aislada}
Para evitar que el encabezado citante induzca emparejamientos espurios con secciones homónimas del citado (dado el alto peso SHAP de los títulos de sección en \textit{SciBERT}~\cite{paolini2025citefusion}), aislamos la consulta al texto discursivo citante (Etapa 1, Figura~\ref{fig:pipeline_dos_etapas}):
\begin{equation}
\label{eq:consulta_aislada}
T_{\text{query}} = x_{\text{ctx}}, \qquad T_{\text{chunk}}(c_{k,j}) = \text{text}(c_{k,j}).
\end{equation}
Con estados ocultos finales $\mathbf{H}^{(L)} = (\mathbf{h}_1^{(L)}, \dots, \mathbf{h}_n^{(L)}) \in \mathbb{R}^{n \times 768}$ y máscara $\mathbf{m} \in \{0, 1\}^n$, aplicamos \textit{mean pooling} enmascarado~\cite{reimers2019sentence} y normalización $L_2$ en $\mathbb{S}^{767}$:
\begin{equation}
\label{eq:mean_pooling_l2}
\mathbf{e}(T) = \frac{\sum_{i=1}^{n} m_i \mathbf{h}_i^{(L)}}{\max\!\left(\sum_{i=1}^{n} m_i, 10^{-9}\right)}, \quad \mathbf{u}(T) = \frac{\mathbf{e}(T)}{\|\mathbf{e}(T)\|_2} \in \mathbb{S}^{767}.
\end{equation}
Los vectores $\mathbf{u}(T)$ se calculan en memoria sin persistencia vectorial en disco, guardando en Capa 4 los $K=3$ fragmentos de mayor similitud coseno $s_{k,j} = \langle \mathbf{u}(T_{\text{query}}), \mathbf{u}(T_{\text{chunk}}(c_{k,j})) \rangle$ ($\bar{s}_{\pi_1} = 0,8943$).

\begin{figure*}[!t]
\centering
\resizebox{0.89\textwidth}{!}{%
\begin{tikzpicture}[
    font=\sffamily\scriptsize,
    module/.style={draw=black!70, rounded corners=2pt, minimum height=6.0mm, inner sep=3pt, align=center, font=\sffamily\scriptsize},
    tokblock/.style={draw=black!65, minimum height=5.5mm, inner sep=2.5pt, align=center, font=\sffamily\scriptsize},
    arrow/.style={-{Latex[length=2.0mm, width=1.4mm]}, semithick, black!80}
]
\node[font=\sffamily\bfseries\scriptsize, anchor=west, text=teal!70!black] at (0, 1.05) {Etapa 1: Recuperación Densa Intra-Documento con SciBERT (Aislada de Encabezados)};
\node[module, fill=blue!10, text width=2.4cm] (qtxt) at (1.3, 0.4) {Consulta $T_{\text{query}} = x_{\text{ctx}}$};
\node[module, fill=teal!10, text width=2.4cm] (ctxt) at (1.3, -0.4) {Chunks $\Omega(\mathcal{P}_{\text{cited}})$ ($\le 110$ pal.)};
\node[module, fill=red!12, text width=2.1cm] (sb1) at (4.3, 0.0) {\textbf{SciBERT} (110M)\\\textit{Mean Pool} $+ L_2$};
\node[module, fill=yellow!18, text width=2.1cm] (cos) at (7.0, 0.0) {Coseno $s_j = \langle \mathbf{u}_q, \mathbf{u}_{c_j} \rangle$};
\node[module, fill=green!12, text width=2.5cm] (top3) at (9.9, 0.0) {\textbf{Top-3 Chunks} $\mathcal{C}_{\text{top3}}$\\$+ \text{Secciones } \tilde{h}_{\pi_r}$};
\draw[arrow] (qtxt.east) -- ([yshift=4pt]sb1.west); \draw[arrow] (ctxt.east) -- ([yshift=-4pt]sb1.west);
\draw[arrow] (sb1.east) -- node[above, font=\tiny] {$\mathbb{S}^{767}$} (cos.west); \draw[arrow] (cos.east) -- (top3.west);

\node[font=\sffamily\bfseries\scriptsize, anchor=west, text=blue!70!black] at (0, -1.15) {Etapa 2: Empaquetado de Pares de Secuencias y Clasificación Supervisada (A4) / Decodificación en $\Delta^8$ (A5)};
\node[tokblock, fill=yellow!25] (cls) at (0.4, -1.78) {\texttt{[CLS]}};
\node[tokblock, fill=violet!15, text width=1.5cm] (secA) at (1.6, -1.78) {Sección $h_{\text{citing}}$};
\node[tokblock, fill=violet!15, text width=2.2cm] (ctxA) at (3.65, -1.78) {Contexto ($m=3$, \texttt{<cite>})};
\node[tokblock, fill=cyan!20] (sep1) at (5.2, -1.78) {\texttt{[SEP]}};
\node[tokblock, fill=green!15, text width=4.5cm] (evB) at (7.9, -1.78) {\textbf{A:} $t_{\text{cited}} + \text{abs}_{\text{cited}}$ $\mid$ \textbf{B:} $\bigoplus_{r=1}^3 (\tilde{h}_{\pi_r}: \tilde{f}_r)$ ($Q_{\text{chunk}}$)};
\node[tokblock, fill=cyan!20] (sep2) at (10.6, -1.78) {\texttt{[SEP]}};
\draw[decorate, decoration={brace, mirror, raise=2pt}, semithick] (cls.south west) -- (ctxA.south east) node[midway, below=3pt, font=\sffamily\tiny] {Segmento 1: Citante ($L_{\text{citing}}$)};
\draw[decorate, decoration={brace, mirror, raise=2pt}, semithick] (evB.south west) -- (sep2.south east) node[midway, below=3pt, font=\sffamily\tiny] {Segmento 2: Evidencia ($512 - L_{\text{citing}} - 3$)};
\draw[arrow, dashed, teal!75!black] (top3.south) -- ++(0,-0.32) -| (evB.north);
\node[module, fill=red!14, text width=2.9cm] (a4mod) at (13.3, -1.35) {\textbf{A4:} \textit{SciBERT}/\textit{LoRA} $\to \text{Softmax}$};
\node[module, fill=orange!15, text width=2.9cm] (a5mod) at (13.3, -2.2) {\textbf{A5:} JSON $+ \text{Proyección } L_1$};
\node[module, fill=yellow!20, text width=1.5cm] (out) at (16.0, -1.78) {\textbf{Símplex $\Delta^8$}\\$\hat{y} \in \mathcal{Y}_9$};
\draw[arrow] (sep2.east) -- ++(0.3,0) |- (a4mod.west); \draw[arrow] (sep2.east) -- ++(0.3,0) |- (a5mod.west);
\draw[arrow] (a4mod.east) -- ([yshift=3pt]out.west); \draw[arrow] (a5mod.east) -- ([yshift=-3pt]out.west);
\end{tikzpicture}%
}
\caption{Pipeline en dos etapas de EvidenCite (inspirado en Gu et al.~\cite{gu2022local}): recuperación densa Top-3 con \textit{SciBERT} (Etapa 1) e inferencia en $\Delta^8$ (Etapa 2). Elaboración propia.}
\label{fig:pipeline_dos_etapas}
\end{figure*}

\subsubsection{Particionamiento Estratificado Agrupado Anti-Fuga}
Para garantizar cero fuga documental ($\text{leakage} = 0$), agrupamos las instancias en bloques disjuntos $G_p = \{i \in \mathcal{I} \mid p_{\text{citing}}(i) = p\}$ por artículo citante~\cite{lauscher2022multicite}. Ordenando los grupos por rareza de su clase más minoritaria, asignamos cada bloque íntegro $G_p$ a la partición $s^* \in \{\text{train}, \text{dev}, \text{test}\}$ ($r_s \in \{0,70; 0,15; 0,15\}$) que minimiza el costo compuesto:
\begin{equation}
\label{eq:split_antifuga}
\begin{aligned}
J(s; G_p) &= \frac{|(N_s + |G_p|) - T_s|}{T_s} \\
&\quad - 2,5 \sum_{c \in \mathcal{Y}(G_p)} \frac{\max\!\left(0, T_{s,c} - (N_{s,c} + n_{G_p,c})\right)}{\max(1, T_{s,c})},
\end{aligned}
\end{equation}
donde $T_s$ y $T_{s,c}$ son las cuotas objetivo en $s$, asegurando $\mathcal{P}_{\text{citing}}(\mathcal{D}_{s_1}) \cap \mathcal{P}_{\text{citing}}(\mathcal{D}_{s_2}) = \emptyset$ ($\forall s_1 \neq s_2$).

\subsection{Homologación Taxonómica, Supervisión Débil NLI, Panel de Jueces y Validación Humana (A2--A3)}

\subsubsection{Homologación Ontológica a las 9 Clases de Lyu et al.}
Proyectamos las 7 etiquetas de \textit{MultiCite} sobre $\mathcal{Y}_9$: \textit{background} $\to$ \texttt{Background}, \textit{uses} $\to$ \texttt{Application}, \textit{extends} $\to$ \texttt{Improvement}, $\textit{similarities} \cup \textit{differences} \to \texttt{Comparison}$, y $\textit{motivation} \cup \textit{future work} \to \texttt{Gap}$ (15 instancias de trabajo futuro). Para cubrir las cuatro clases ausentes en \textit{MultiCite} (\texttt{Basis}, \texttt{Evidence}, \texttt{Originator}, \texttt{Further Reading})~\cite{lyu2021classification}, combinamos $4.445$ instancias minadas con patrones léxicos de alta especificidad en \textit{unarXive} y supervisión débil NLI.

\subsubsection{Teorema de Co-Localización y Calibración NLI en \texorpdfstring{$\Delta^8$}{Delta-8}}
Evaluamos el pool no etiquetado con \textit{DeBERTa-v3-base}~\cite{he2023debertav3} (86M parámetros, MNLI/FEVER/ANLI). Cada premisa incluye el encabezado como marco ($P_i = \text{``Section: ''} \oplus h_{\text{citing}} \oplus \text{`` . Context: ''} \oplus x_{\text{ctx}}$) frente a 9 hipótesis canónicas $\{H_k\}_{k=1}^9$, proyectando los logits $(e_{i,1}, \dots, e_{i,9})$ sobre $\Delta^8$ con temperatura $\tau$:
\begin{equation}
\label{eq:nli_softmax_temp}
P(c_k \mid P_i) = \frac{\exp(e_{i,k} / \tau)}{\sum_{j=1}^{9} \exp(e_{i,j} / \tau)}.
\end{equation}
En una prueba preliminar sobre $11.000$ candidatos con umbrales ingenuos ($\tau = 0,7, \theta_{\text{conf}} = 0,70, \theta_{\text{margin}} = 0,15$), las clases unimodales (\texttt{Further Reading}, \texttt{Evidence}) saturaron sus cuotas mientras la tríada metodológica (\texttt{Basis}, \texttt{Improvement}, \texttt{Comparison}) sufrió rechazo masivo. Ello ocurre porque en ventanas de $m=3$ oraciones el discurso suele combinar adopción de base (\texttt{Basis}), modificación técnica (\texttt{Improvement}) y contraste empírico (\texttt{Comparison})~\cite{jurgens2018measuring,pride2020authoritative}. Sean $P_1 \ge P_2$ las probabilidades posteriores de las dos funciones co-activas, $\epsilon = \sum_{k=3}^9 P_{(k)} \approx 0,19$ la masa residual de las 7 clases inactivas y $\Delta = P_1 - P_2$ el margen; por conservación en $\Delta^8$ ($P_1 + P_2 = 1 - \epsilon$):
\begin{equation}
\label{eq:teorema_colocalizacion}
P_1 = \frac{1 - \epsilon + \Delta}{2}, \qquad P_2 = \frac{1 - \epsilon - \Delta}{2}.
\end{equation}
Con un margen amplio $\Delta = 0,35$ y $\epsilon = 0,19$, el máximo teórico es $P_1 = (0,81 + 0,35)/2 = 0,580 < 0,70$; exigir $\theta_{\text{conf}} = 0,70$ imponía $\Delta \ge 0,59$, excluyendo toda co-ocurrencia legítima. Calibramos sobre $N=1.000$ candidatos el régimen $(\tau = 0,6, \theta_{\text{conf}} = 0,50, \theta_{\text{margin}} = 0,15)$ ($P_1 \ge 4,5\times$ el azar $1/9$), elevando la admisión en $2,44\times$ ($6,6\pct{} \to 16,1\pct{}$) y multiplicando en $5,5\times\text{--}13\times$ la captura de \texttt{Basis} e \texttt{Improvement}.

\subsubsection{Panel de 3 Jueces Abiertos en vLLM y Regla Conservadora (A2)}
Para auditar las $15.029$ citas minadas y medir calibración sobre las $3.629$ Gold de \textit{MultiCite} sin contaminar los modelos de A5, desplegamos en \textit{vLLM}~\cite{kwon2023efficient} (NVIDIA A100-40GB, $\text{temp.}=0,0$) tres jueces abiertos independientes: \texttt{Llama-3.1-8B-Instruct}, \texttt{Qwen2.5-7B-Instruct} y \texttt{Mistral-7B-Instruct-v0.3}, bajo \textbf{cegamiento del citado} y sin sección citante ($93,6\pct{}$ ausente en \textit{MultiCite}). Cada juez emite justificación ($\le 300$ car.) y voto $v_j \in \mathcal{Y}_9$. Con frecuencia modal $c_{\text{top}} \in \{1, 2, 3\}$, moda $L_{\text{maj}}$ y etiqueta previa $y_{\text{prior}}$:
\begin{equation}
\label{eq:reglas_consenso_panel}
\resizebox{0.87\columnwidth}{!}{$\displaystyle
(y_{\text{A2}}, \ell_{\text{src}}) = \begin{cases}
(y_{\text{prior}}, \text{Gold}), & \text{Caso F: } b_{\text{gold}}, \\
(L_{\text{maj}}, \text{Panel}), & \text{Caso A: } \neg b_{\text{gold}} \land c_{\text{top}} = 3, \\
(y_{\text{prior}}, \text{Panel}), & \text{Caso B: } \neg b_{\text{gold}} \land c_{\text{top}} = 2 \land L_{\text{maj}} = y_{\text{prior}}, \\
(y_{\text{prior}}, \text{Rev.}), & \text{Caso C: } \neg b_{\text{gold}} \land c_{\text{top}} = 2 \land L_{\text{maj}} \neq y_{\text{prior}}, \\
(y_{\text{prior}}, \text{Rev.}), & \text{Caso D: } \neg b_{\text{gold}} \land c_{\text{top}} = 1.
\end{cases}
$}
\end{equation}
Retener $y_{\text{prior}}$ en el Caso C evita que el sesgo del panel 7B--8B hacia las tres clases genéricas (\texttt{Comparison}, \texttt{Application}, \texttt{Background}: $65\pct{}\text{--}68\pct{}$ de sus votos) erosione en $39\pct{}\text{--}53\pct{}$ las clases minoritarias (\texttt{Originator}, \texttt{Evidence}, \texttt{Improvement}).

\subsubsection{Validación y Curaduría Humana Experta del Conjunto de Prueba (A3)}
En la Actividad A3, expertos humanos revisaron manualmente el $100\pct{}$ de $\mathcal{D}_{\text{test}}$ ($N = 2.798$ instancias, $1.258$ citantes), certificando cada registro como verdad terreno humana ($b_{\text{gold}} = \text{True}$) y conservando $\mathcal{D}_{\text{test}}^{(\text{panel})}$ para auditar el cambio distribucional.

\subsection{Clasificadores Supervisados y Adaptadores PEFT/LoRA (A4)}

La Actividad A4 entrena tres arquitecturas discriminativas sobre $\mathcal{D}_{\text{train}}$ ($N=13.077$), validadas en $\mathcal{D}_{\text{dev}}$ ($N=2.783$):

\subsubsection{Línea Base Léxica (TF-IDF + Regresión Logística)}
Extrae unigramas y bigramas ($\text{rango} = (1, 2)$, frecuencia documental $[2, 0,95]$, $\le 50.000$ atributos) con $\text{tf}'(t, d) = 1 + \log(\text{tf}(t, d))$ y Regresión Logística multinomial $L_2$ ($C = 1,0$) con pesos balanceados $w_c = N / (9 \cdot n_c)$.

\subsubsection{Ajuste Fino de SciBERT con Presupuesto Equitativo por Fragmento}
En \textit{SciBERT} (110M, $L=12, d=768$)~\cite{beltagy2019scibert}, empaquetamos pares de secuencias estilo codificador cruzado~\cite{gu2022local} (Etapa 2, Figura~\ref{fig:pipeline_dos_etapas}):
\begin{equation}
\label{eq:empaquetado_scibert}
\mathbf{x}_{\text{input}} = \texttt{[CLS]} \;\; T_{\text{citing}} \;\; \texttt{[SEP]} \;\; T_{\text{cited}}^{(\mu)} \;\; \texttt{[SEP]}.
\end{equation}
En \textbf{Config A}, $T_{\text{cited}}^{(A)} = t_{\text{cited}} \oplus_{\texttt{". "}} \text{abs}_{\text{cited}}$. En \textbf{Config B}, cada pasaje antepone su sección autoral $f_r = \tilde{h}_{\pi_r} \oplus \texttt{": "} \oplus \text{text}(c_{\pi_r})$. Reservando $L_{\text{reserved}} = L_{\text{citing}} + 3$, asignamos a cada fragmento una cuota equitativa $Q_{\text{chunk}}$ dentro de $L_{\max} = 512$:
\begin{equation}
\label{eq:cuota_equitativa_chunks}
Q_{\text{chunk}} = \max\!\left(1, \left\lfloor \frac{\max(K, L_{\max} - L_{\text{citing}} - 3) - K + 1}{K} \right\rfloor\!\right)\!.
\end{equation}
Recortando cada $f_r$ a $Q_{\text{chunk}}$ sub-palabras ($T_{\text{cited}}^{(B)} = \tilde{f}_1 \oplus_{\texttt{". "}} \tilde{f}_2 \oplus_{\texttt{". "}} \tilde{f}_3$), proyectamos $\mathbf{h}_{\texttt{[CLS]}} \in \mathbb{R}^{768}$ (\textit{dropout} $p=0,1$) hacia $\mathbf{z} \in \mathbb{R}^9$ y optimizamos con \textit{AdamW} ($\eta = 2\times 10^{-5}$, $\lambda = 0,01$, 5 épocas, \textit{warmup} $10\pct{}$) bajo entropía cruzada ponderada:
\begin{equation}
\label{eq:weighted_ce_loss}
\mathcal{L}_{\text{WCE}}\!\left(\{(\mathbf{z}_i, y_i)\}_{i=1}^{B}\right) = \frac{\sum_{i=1}^{B} w_{y_i} \left(-\log \text{Softmax}(\mathbf{z}_i)_{y_i}\right)}{\sum_{i=1}^{B} w_{y_i}}.
\end{equation}

\subsubsection{Adaptación de Bajo Rango (PEFT/LoRA) sobre Modelos Abiertos}
En decodificadores abiertos (\texttt{Qwen-8B}), congelamos $\mathbf{W}_0 \in \mathbb{R}^{d_{\text{out}} \times d_{\text{in}}}$ e inyectamos adaptadores $\mathbf{A} \in \mathbb{R}^{r \times d_{\text{in}}}, \mathbf{B} \in \mathbb{R}^{d_{\text{out}} \times r}$ ($r = 16, \alpha = 32$)~\cite{hu2022lora} entrenados en 32 bits ($\eta = 10^{-4}$, lote $16$):
\begin{equation}
\label{eq:lora_update}
\mathbf{h} = \mathbf{W}_0 \mathbf{x} + \frac{\alpha}{r} \, \mathbf{B}\mathbf{A}\mathbf{x}.
\end{equation}

\subsection{Evaluación de LLMs Comerciales de Frontera vía API (A5)}

\subsubsection{Diseño Factorial \texorpdfstring{$2 \times 2 \times 2$}{2x2x2} y Control de Razonamiento Latente}
Evaluamos sobre $\mathcal{D}_{\text{test}}$ ($N = 2.798$ Gold) un diseño factorial $2 \times 2 \times 2$ ($8$ variantes, $22.384$ inferencias): (1) \textbf{Proveedor:} \texttt{OpenAI gpt-4.1-mini} ($\text{temp.}=0,0$, semilla $42$, $0$ tokens de pensamiento, salida $\le 512$ tokens; $\$0,40/1\text{M}$ entrada, $\$0,10/1\text{M}$ caché, $\$1,60/1\text{M}$ salida) vs. \texttt{Google gemini-3.8-flash} en Vertex AI (pensamiento mínimo \texttt{LOW}, semilla $42$, salida $\le 2.048$ tokens; $\$0,75/1\text{M}$ entrada, $\$0,075/1\text{M}$ caché, $\$3,75/1\text{M}$ salida+pensamiento); (2) \textbf{Evidencia ($\text{Config A}$ vs. $\text{B}$):} sin sección citante en usuario por simetría con \textit{MultiCite}, pero preservando en Config B los encabezados autorales $\tilde{h}_{\pi_r}$ de los Top-3 fragmentos; y (3) \textbf{Régimen ($\text{0-Shot}$ vs. $\text{9-Shot}$):} $k=0$ vs. $k=9$ ejemplares (1 por clase, semilla $42$) muestreados de $\mathcal{D}_{\text{train}}$ con paridad 1:1.

\subsubsection{Decodificación Restringida a Esquema Estructurado y Confianza Verbalizada vs. Softmax Paramétrico}
Imponemos una gramática JSON estricta (\textit{Structured Outputs}) con etiqueta $\hat{y} \in \mathcal{Y}_9$ y diccionario de nueve puntajes no negativos $\mathbf{s} = (s_1, \dots, s_9)^\top \in [0, 1]^9$ (sin justificación textual), proyectados post-hoc mediante normalización $L_1$ sobre $\Delta^8$:
\begin{equation}
\label{eq:proyeccion_l1_simplex}
p_k = \begin{cases}
s_k / \sum_{j=1}^{9} s_j, & \text{si } \sum_{j=1}^{9} s_j > 0, \\[4pt]
1 / 9, & \text{si } \sum_{j=1}^{9} s_j = 0 \text{ (respaldo)}.
\end{cases}
\end{equation}
A diferencia de A4 —donde $\mathbf{p} = \text{Softmax}(\mathbf{W}\mathbf{h} + \mathbf{b})$ proviene de logits paramétricos—, en A5 $\mathbf{p}$ constituye \textbf{confianza verbalizada estructurada}: dígitos emitidos autorregresivamente y proyectados vía Ec.~\eqref{eq:proyeccion_l1_simplex}, no $\exp(\text{logprobs})$ de tokens internos.

\subsubsection{Construcción Determinista de Ejemplares 9-Shot Guiada por Vecindad Retórica en \texorpdfstring{$\Delta^8$}{Delta-8}}
Para evitar que vectores fijos (p.~ej., $0,92$ y ocho $0,01$) induzcan copia de plantilla (\textit{template copying}), particionamos para cada clase $c^* \in \mathcal{Y}_9$ las 8 clases restantes según la auditoría de Lyu et al.~\cite{lyu2021classification} y la Ec.~\eqref{eq:teorema_colocalizacion} en dos \textbf{vecinos primarios} $\mathcal{N}_1(c^*)$, dos \textbf{secundarios} $\mathcal{N}_2(c^*)$ y cuatro de \textbf{cola ortogonal} $\mathcal{N}_{\perp}(c^*)$, con banda dominante $[L_{c^*}, U_{c^*}]$ (Tabla~\ref{tab:vecindad_retorica}).

\begin{table}[!t]
\centering
\caption{Matriz de Vecindad Retórica en $\Delta^8$ para Síntesis Determinista de Distribuciones en Demostraciones 9-Shot}
\label{tab:vecindad_retorica}
\scriptsize
\setlength{\tabcolsep}{2.4pt}
\renewcommand{\arraystretch}{1.04}
\resizebox{\columnwidth}{!}{%
\begin{tabular}{l c l l}
\toprule
\textbf{Clase Real ($c^*$)} & \textbf{Banda $[L_{c^*}, U_{c^*}]$} & \textbf{Vecinos Primarios $\mathcal{N}_1(c^*)$} & \textbf{Vecinos Secundarios $\mathcal{N}_2(c^*)$} \\
\midrule
\texttt{Background} & $[0,62,\; 0,78]$ & \texttt{Further Reading}, \texttt{Gap} & \texttt{Originator}, \texttt{Evidence} \\
\texttt{Gap} & $[0,60,\; 0,76]$ & \texttt{Background}, \texttt{Comparison} & \texttt{Evidence}, \texttt{Improvement} \\
\texttt{Basis} & $[0,56,\; 0,72]$ & \texttt{Improvement}, \texttt{Application} & \texttt{Background}, \texttt{Originator} \\
\texttt{Comparison} & $[0,58,\; 0,75]$ & \texttt{Improvement}, \texttt{Evidence} & \texttt{Gap}, \texttt{Background} \\
\texttt{Application} & $[0,60,\; 0,76]$ & \texttt{Basis}, \texttt{Improvement} & \texttt{Evidence}, \texttt{Comparison} \\
\texttt{Improvement} & $[0,56,\; 0,72]$ & \texttt{Basis}, \texttt{Comparison} & \texttt{Application}, \texttt{Evidence} \\
\texttt{Evidence} & $[0,64,\; 0,80]$ & \texttt{Comparison}, \texttt{Background} & \texttt{Basis}, \texttt{Gap} \\
\texttt{Originator} & $[0,64,\; 0,80]$ & \texttt{Background}, \texttt{Basis} & \texttt{Application}, \texttt{Evidence} \\
\texttt{Further Reading} & $[0,68,\; 0,84]$ & \texttt{Background}, \texttt{Evidence} & \texttt{Application}, \texttt{Basis} \\
\bottomrule
\end{tabular}%
}
\end{table}

Con semilla de 64 bits derivada de $\text{SHA-256}(id_{\text{ex}} \parallel \texttt{"::"} \parallel c^*)$, muestreamos $g \in [100 L_{c^*}, 100 U_{c^*}]$ (en centésimas) y repartimos el remanente $R = 100 - g$ de forma estratificada:
\begin{equation}
\label{eq:reparto_vecindad_9shot}
\begin{aligned}
B_1 &= \text{round}(R \cdot u_1), \quad u_1 \sim \mathcal{U}(0,58; 0,70) \;\; (\text{hacia } \mathcal{N}_1(c^*)), \\
B_2 &= \text{round}(R \cdot u_2), \quad u_2 \sim \mathcal{U}(0,20; 0,28) \;\; (\text{hacia } \mathcal{N}_2(c^*)), \\
B_\perp &= R - B_1 - B_2 \;\; (\text{reparto discreto en } \mathcal{N}_\perp(c^*)),
\end{aligned}
\end{equation}
con proporciones internas $\eta_1 \sim \mathcal{U}(0,55; 0,72)$ y $\eta_2 \sim \mathcal{U}(0,52; 0,75)$, asegurando $\sum_{k=1}^9 s_k = 1,00$ y $s_k < s_{c^*}$ ($\forall k \neq c^*$).

\subsubsection{Contabilidad Financiera Dual (Bruta vs. Neta) bajo Prompt Caching}
En \textit{9-Shot}, el prefijo invariante de las $N = 2.798$ llamadas ($\approx 3.968$ tokens en Config A y $\approx 6.016$ en Config B) activa el caché de atención Key-Value en servidor (\textit{Prompt Caching}). Registrando por instancia $i$ los tokens de entrada $I_i$, caché leído $C_i \le I_i$, salida $O_i$ y pensamiento $T_i$, computamos el \textbf{costo bruto teórico} ($K_{\text{gross}}$, tarifa plena $\pi_{\text{in}}$) y el \textbf{costo neto real} ($K_{\text{net}}$, tarifa con descuento $\pi_{\text{cache}}$):
\begin{equation}
\label{eq:costos_bruto_neto}
\begin{aligned}
K_{\text{gross}} &= \frac{1}{10^6} \sum_{i=1}^{N} \left[ I_i \, \pi_{\text{in}} + (O_i + T_i) \, \pi_{\text{out}} \right], \\
K_{\text{net}} &= \frac{1}{10^6} \sum_{i=1}^{N} \left[ (I_i - C_i) \, \pi_{\text{in}} + C_i \, \pi_{\text{cache}} + (O_i + T_i) \, \pi_{\text{out}} \right].
\end{aligned}
\end{equation}

\subsection{Arquitectura de Despliegue Cloud y Demostrador Interactivo (A8)}

La Actividad A8 despliega la inferencia al vuelo en cuatro etapas: (1) \textbf{Normalización Citante:} protege la primera ocurrencia de la cita objetivo $c_{\text{target}}$, enmascara las co-citas restantes como \texttt{[OTHERCIT]} y restituye \texttt{<cite>}$c_{\text{target}}$\texttt{</cite>}; (2) \textbf{Resolución Jerárquica L1--L4:} consulta en cascada (L1) cachés RAM de predictores, embeddings \textit{SciBERT} y PDFs por hash SHA-256, (L2) almacén local, (L3) API de arXiv, HTML \texttt{ar5iv} y Semantic Scholar, y (L4) extracción estructural en memoria de PDFs ($\le 20\text{ MB}$, $\le 20$ pág.) en $\sim 3\text{--}5\text{ s}$; (3) \textbf{Recuperación Top-3 en Tiempo Real:} segmenta con la Ec.~\eqref{eq:chunking_voraz} ($\le 110$ pal.) y recupera con \textit{SciBERT} vía Ec.~\eqref{eq:mean_pooling_l2}; y (4) \textbf{Despacho Cloud:} un patrón estrategia unifica clasificadores locales y APIs comerciales tras un servicio REST (\textit{FastAPI}) e interfaz web (\textit{Nginx}) contenerizados en Docker y desplegados en \textbf{AWS ECS Fargate} detrás de un ALB aprovisionado con \textit{Terraform} ($2\text{ vCPU} / 8\text{ GB RAM}$ para API y $0,25\text{ vCPU} / 512\text{ MB}$ para cliente web).

\section{Resultados Experimentales}

\subsection{Consolidación del Corpus 1:1, Auditoría del Panel A2 y Validación Humana A3}

\subsubsection{Métricas Estructurales de A1 y Techo Físico de Clases Minoritarias}
El pipeline A1 ($61.100,16\text{ s}$) consolidó $\mathbf{18.658}$ instancias con paridad estricta 1:1 ($18.658$ para Config A y $18.658$ para Config B) y fuga cero entre particiones ($\text{leakage} = 0$: $\mathcal{D}_{\text{train}}=13.077$ [$70,09\pct{}$, $5.835$ citantes], $\mathcal{D}_{\text{dev}}=2.783$ [$14,92\pct{}$, $1.243$ citantes] y $\mathcal{D}_{\text{test}}=2.798$ [$15,00\pct{}$, $1.258$ citantes]). La recuperación Top-3 con \textit{SciBERT} logró una similitud Top-1 media de $\mathbf{0,8943}$ sobre ventanas de $\mathbf{83,94}$ palabras. En supervisión débil NLI, \textit{DeBERTa-v3-base} evaluó $135.035$ candidatos de \textit{unarXive}, aceptando $\mathbf{63.802}$ ($\mathbf{47,25\pct{}}$) y rechazando $71.233$ ($70.461$ por $\theta_{\text{conf}} < 0,50$ y $772$ por $\theta_{\text{margin}} < 0,15$). Siete clases alcanzaron el tope con búfer ($2.300$ en Capa 4, Tabla~\ref{tab:distribucion_etapas_a1_a2_a3}), mientras \texttt{Application} ($2.028$), \texttt{Improvement} ($1.641$) y \texttt{Basis} ($1.189$) reflejan su techo físico bajo texto completo citado en arXiv.

\begin{table*}[!t]
\centering
\caption{Evolución Auditada del Corpus EvidenCite ($N=18.658$) a través de las Etapas A1 (Capa 4), A2 (Panel de 3 Jueces) y A3 (Curaduría Humana Gold de $\mathcal{D}_{\text{test}}$, $N=2.798$)}
\label{tab:distribucion_etapas_a1_a2_a3}
\scriptsize
\setlength{\tabcolsep}{4.0pt}
\renewcommand{\arraystretch}{1.04}
\resizebox{\textwidth}{!}{%
\begin{tabular}{l r r r r r r r r r}
\toprule
& \multicolumn{2}{c}{\textbf{Etapa A1 (Capa 4 Pre-Panel)}} & \multicolumn{3}{c}{\textbf{Etapa A2 (Post-Panel 3 Jueces vLLM)}} & \multicolumn{4}{c}{\textbf{Etapa A3 (Corpus Final con Test 100\pct{} Gold Humano)}} \\
\cmidrule(lr){2-3} \cmidrule(lr){4-6} \cmidrule(lr){7-10}
\textbf{Clase Canónica ($\mathcal{Y}_9$)} & \textbf{Aceptados NLI} & \textbf{Total A1} & \textbf{Total A2} & \textbf{Déficit ($2\text{k}$)} & \textbf{Test Pre-A3} & \textbf{Train Final} & \textbf{Dev Final} & \textbf{Test Gold A3 ($\pct{}$)} & \textbf{Total Final} \\
\midrule
\texttt{Background} & $7.115$ & $2.300$ & $2.369$ & $0$ & $348$ ($12,4\%$) & $1.671$ & $350$ & $\mathbf{1.200}$ ($\mathbf{42,89\%}$) & $\mathbf{3.221}$ \\
\texttt{Application} & $857$ & $2.028$ & $2.838$ & $0$ & $453$ ($16,2\%$) & $1.974$ & $411$ & $\mathbf{637}$ ($\mathbf{22,77\%}$) & $\mathbf{3.022}$ \\
\texttt{Comparison} & $1.618$ & $2.300$ & $3.053$ & $0$ & $493$ ($17,6\%$) & $2.119$ & $441$ & $\mathbf{298}$ ($\mathbf{10,65\%}$) & $\mathbf{2.858}$ \\
\texttt{Further Reading} & $13.611$ & $2.300$ & $2.001$ & $0$ & $251$ ($9,0\%$) & $1.453$ & $297$ & $\mathbf{227}$ ($\mathbf{8,11\%}$) & $\mathbf{1.977}$ \\
\texttt{Gap} & $3.146$ & $2.300$ & $2.181$ & $0$ & $327$ ($11,7\%$) & $1.511$ & $343$ & $\mathbf{73}$ ($\mathbf{2,61\%}$) & $\mathbf{1.927}$ \\
\texttt{Originator} & $3.937$ & $2.300$ & $1.751$ & $249$ & $257$ ($9,2\%$) & $1.228$ & $266$ & $\mathbf{105}$ ($\mathbf{3,75\%}$) & $\mathbf{1.599}$ \\
\texttt{Evidence} & $31.408$ & $2.300$ & $1.716$ & $284$ & $253$ ($9,0\%$) & $1.195$ & $268$ & $\mathbf{100}$ ($\mathbf{3,57\%}$) & $\mathbf{1.563}$ \\
\texttt{Improvement} & $1.401$ & $1.641$ & $1.421$ & $579$ & $217$ ($7,8\%$) & $998$ & $206$ & $\mathbf{68}$ ($\mathbf{2,43\%}$) & $\mathbf{1.272}$ \\
\texttt{Basis} & $709$ & $1.189$ & $1.328$ & $672$ & $199$ ($7,1\%$) & $928$ & $201$ & $\mathbf{90}$ ($\mathbf{3,22\%}$) & $\mathbf{1.219}$ \\
\midrule
\textbf{Total Instancias} & $\mathbf{63.802}$ & $\mathbf{18.658}$ & $\mathbf{18.658}$ & $\mathbf{1.784}$ & $\mathbf{2.798}$ ($100\%$) & $\mathbf{13.077}$ & $\mathbf{2.783}$ & $\mathbf{2.798}$ ($\mathbf{100,00\%}$) & $\mathbf{18.658}$ \\
\bottomrule
\end{tabular}%
}
\end{table*}

\subsubsection{Auditoría del Panel de 3 Jueces Abiertos (A2)}
En A2, los tres jueces en \textit{vLLM} evaluaron las $18.658$ instancias con abstención casi nula ($0\pct{}$ en \texttt{Llama-3.1-8B} [$10,20\text{M}$ tokens entrada, $0,88\text{M}$ salida, $1.600,09\text{ ms}$] y \texttt{Qwen2.5-7B} [$9,88\text{M}$ entrada, $0,72\text{M}$ salida, $1.311,28\text{ ms}$]; $23$ abstenciones [$0,123\pct{}$] en \texttt{Mistral-7B} [$11,14\text{M}$ entrada, $1,14\text{M}$ salida, $1.363,17\text{ ms}$]). Bajo la Ec.~\eqref{eq:reglas_consenso_panel}, resultaron: \textbf{Caso A} (3/3) $= \mathbf{4.772}$ ($25,58\pct{}$); \textbf{Caso B} (2/3 confirmatorio) $= \mathbf{2.667}$ ($14,29\pct{}$); \textbf{Caso C} (2/3 divergente retenido) $= \mathbf{4.705}$ ($25,22\pct{}$); \textbf{Caso D} (1-1-1) $= \mathbf{2.885}$ ($15,46\pct{}$); y \textbf{Caso F} (Gold \textit{MultiCite}) $= \mathbf{3.629}$ ($19,45\pct{}$), sumando $3.629$ Gold, $7.439$ consensuadas ($A+B$: $5.389 / 10.584 = 50,92\pct{}$ en NLI y $2.050 / 4.445 = 46,12\pct{}$ en léxicas) y $7.590$ en revisión ($C+D$). El $\kappa$ de Cohen por pares fue $\mathbf{0,4388}$ (\texttt{Llama}--\texttt{Qwen}), $\mathbf{0,3881}$ (\texttt{Mistral}--\texttt{Qwen}) y $\mathbf{0,3548}$ (\texttt{Llama}--\texttt{Mistral}), con media $\bar{\kappa} = \mathbf{0,3939}$. Sobre las $3.629$ Gold de \textit{MultiCite}, el voto mayoritario alcanzó $\mathbf{36,40\pct{}}$ (\texttt{Qwen}: $36,92\pct{}$, \texttt{Llama}: $35,27\pct{}$, \texttt{Mistral}: $31,36\pct{}$), desviando \texttt{Background} hacia \texttt{Comparison} ($410$), \texttt{Gap} ($179$) y \texttt{Basis} ($173$)~\cite{fogelson2025llms}.

\subsubsection{Impacto de la Curaduría Humana Experta en A3 sobre el Conjunto de Prueba}
La revisión humana del $100\pct{}$ de $\mathcal{D}_{\text{test}}$ ($N = 2.798$) en A3 corrigió $\mathbf{1.520}$ etiquetas ($\mathbf{54,32\pct{}}$) respecto a $\mathcal{D}_{\text{test}}^{(\text{panel})}$: el $69,73\pct{}$ ($781/1.120$) en \textit{Needs Review}, el $48,25\pct{}$ ($538/1.115$) en consenso del panel y el $35,70\pct{}$ ($201/563$) en \textit{MultiCite} (desambiguación a 9 clases). La Tabla~\ref{tab:distribucion_etapas_a1_a2_a3} muestra cómo A3 revirtió la inflación por cuotas en \texttt{Gap} ($327 \to 73$) y \texttt{Comparison} ($493 \to 298$), restituyendo la distribución natural del discurso~\cite{jurgens2018measuring,zhang2025deep} dominada por \texttt{Background} ($348 \to \mathbf{1.200}$, $42,89\pct{}$) y \texttt{Application} ($453 \to \mathbf{637}$, $22,77\pct{}$), y elevando el total Gold humano del corpus a $\mathbf{5.864}$ ($3.066$ en $\mathcal{D}_{\text{train}} \cup \mathcal{D}_{\text{dev}}$ y $2.798$ en $\mathcal{D}_{\text{test}}$).

\subsection{Desempeño Predictivo, Contraste Config A vs. B y Calibración Probabilística (A4, A5, A6)}

\subsubsection{Desempeño Global y Liderazgo de Gemini-3.8-Flash · 0-Shot · Config B}
La Tabla~\ref{tab:maestra_resultados} resume el desempeño sobre el Gold humano ($N = 2.798$). Las filas supervisadas de A4 registran \texttt{[PENDIENTE\_A4]} hasta concluir su entrenamiento en GPU sobre la partición final, analizando aquí las 8 variantes reales de A5. \textbf{\texttt{Gemini-3.8-Flash · 0-Shot · Config B}} lidera globalmente en las tres métricas primarias: $\textbf{Macro-F1} = \mathbf{0,5996}$, $\textbf{Weighted-F1} = \mathbf{0,6933}$ y $\textbf{Accuracy} = \mathbf{0,6841}$ ($68,41\pct{}$), con el menor Brier Score ($\mathbf{0,4809}$). En OpenAI también lidera la Configuración B (\textbf{\texttt{GPT-4.1-mini · 0-Shot · Config B}}: $\text{Macro-F1} = 0,4664$, $\text{Weighted-F1} = 0,5642$, $\text{Accuracy} = 0,5307$).

\begin{table*}[!t]
\centering
\caption{Tabla Maestra Comparativa de Desempeño Predictivo (A6), Calibración Probabilística y Costo Operativo (A7) sobre el Conjunto de Prueba Gold Revisado Manualmente ($N = 2.798$)}
\label{tab:maestra_resultados}
\scriptsize
\setlength{\tabcolsep}{2.8pt}
\renewcommand{\arraystretch}{1.04}
\resizebox{\textwidth}{!}{%
\begin{tabular}{l l c c c c c c c c c}
\toprule
\textbf{Familia / Modelo Evaluado} & \textbf{Régimen} & \textbf{Config.} & \textbf{Macro-F1} $\uparrow$ & \textbf{Weighted-F1} $\uparrow$ & \textbf{Accuracy} $\uparrow$ & \textbf{ECE (10b)} $\downarrow$ & \textbf{Brier} $\downarrow$ & \textbf{NLL} $\downarrow$ & \textbf{Lat. p50 (ms)} $\downarrow$ & \textbf{Costo Neto / 10k} $\downarrow$ \\
\midrule
TF-IDF + Regresión Logística (A4) & Supervisado & A & \texttt{[PENDIENTE\_A4]} & \texttt{[PENDIENTE\_A4]} & \texttt{[PENDIENTE\_A4]} & \texttt{[PENDIENTE\_A4]} & \texttt{[PENDIENTE\_A4]} & \texttt{[PENDIENTE\_A4]} & \texttt{[PENDIENTE\_A4]} & $\$0,00$ (Local) \\
TF-IDF + Regresión Logística (A4) & Supervisado & B & \texttt{[PENDIENTE\_A4]} & \texttt{[PENDIENTE\_A4]} & \texttt{[PENDIENTE\_A4]} & \texttt{[PENDIENTE\_A4]} & \texttt{[PENDIENTE\_A4]} & \texttt{[PENDIENTE\_A4]} & \texttt{[PENDIENTE\_A4]} & $\$0,00$ (Local) \\
\textit{SciBERT} Fine-Tuned 110M (A4) & Supervisado & A & \texttt{[PENDIENTE\_A4]} & \texttt{[PENDIENTE\_A4]} & \texttt{[PENDIENTE\_A4]} & \texttt{[PENDIENTE\_A4]} & \texttt{[PENDIENTE\_A4]} & \texttt{[PENDIENTE\_A4]} & \texttt{[PENDIENTE\_A4]} & $\$0,00$ (Local) \\
\textit{SciBERT} Fine-Tuned 110M (A4) & Supervisado & B & \texttt{[PENDIENTE\_A4]} & \texttt{[PENDIENTE\_A4]} & \texttt{[PENDIENTE\_A4]} & \texttt{[PENDIENTE\_A4]} & \texttt{[PENDIENTE\_A4]} & \texttt{[PENDIENTE\_A4]} & \texttt{[PENDIENTE\_A4]} & $\$0,00$ (Local) \\
LLM Abierto + PEFT \textit{LoRA} (A4) & Supervisado & A & \texttt{[PENDIENTE\_A4]} & \texttt{[PENDIENTE\_A4]} & \texttt{[PENDIENTE\_A4]} & \texttt{[PENDIENTE\_A4]} & \texttt{[PENDIENTE\_A4]} & \texttt{[PENDIENTE\_A4]} & \texttt{[PENDIENTE\_A4]} & $\$0,00$ (Local) \\
LLM Abierto + PEFT \textit{LoRA} (A4) & Supervisado & B & \texttt{[PENDIENTE\_A4]} & \texttt{[PENDIENTE\_A4]} & \texttt{[PENDIENTE\_A4]} & \texttt{[PENDIENTE\_A4]} & \texttt{[PENDIENTE\_A4]} & \texttt{[PENDIENTE\_A4]} & \texttt{[PENDIENTE\_A4]} & $\$0,00$ (Local) \\
\midrule
\texttt{GPT-4.1-mini} (A5) & 0-Shot ($k=0$) & A & $0,4493$ & $0,5374$ & $0,5021$ & $0,3284$ & $0,8449$ & $5,6807$ & $\mathbf{915,9}$ & $\mathbf{\$5,83}$ \\
\texttt{GPT-4.1-mini} (A5) & 0-Shot ($k=0$) & B & $0,4664$ & $0,5642$ & $0,5307$ & $0,3360$ & $0,8221$ & $6,0666$ & $1.025,2$ & $\$6,59$ \\
\texttt{GPT-4.1-mini} (A5) & 9-Shot ($k=9$) & A & $0,4146$ & $0,4960$ & $0,4643$ & $0,3282$ & $0,8531$ & $2,0473$ & $986,5$ & $\$8,23$ \\
\texttt{GPT-4.1-mini} (A5) & 9-Shot ($k=9$) & B & $0,4242$ & $0,5355$ & $0,4964$ & $0,3036$ & $0,8188$ & $2,0760$ & $1.045,6$ & $\$11,61$ \\
\midrule
\texttt{Gemini-3.8-Flash} (A5) & 0-Shot ($k=0$) & A & $0,5900$ & $0,6900$ & $0,6798$ & $0,1204$ & $0,4834$ & $1,2253$ & $2.354,0$ & $\$16,55$ \\
\textbf{\texttt{Gemini-3.8-Flash} (A5)} & \textbf{0-Shot ($k=0$)} & \textbf{B} & $\mathbf{0,5996}$ & $\mathbf{0,6933}$ & $\mathbf{0,6841}$ & $0,1158$ & $\mathbf{0,4809}$ & $1,2778$ & $2.324,6$ & $\$17,68$ \\
\texttt{Gemini-3.8-Flash} (A5) & 9-Shot ($k=9$) & A & $0,5974$ & $0,6693$ & $0,6565$ & $\mathbf{0,1139}$ & $0,5012$ & $\mathbf{1,0687}$ & $2.292,4$ & $\$39,43$ \\
\texttt{Gemini-3.8-Flash} (A5) & 9-Shot ($k=9$) & B & $0,5943$ & $0,6652$ & $0,6523$ & $0,1257$ & $0,5114$ & $1,0913$ & $2.481,1$ & $\$53,09$ \\
\bottomrule
\end{tabular}%
}
\end{table*}

\subsubsection{Desglose por Clase Canónica, Matrices de Confusión y Sección Citante}
Según la Tabla~\ref{tab:f1_por_clase_a5} y la Figura~\ref{fig:desempeno_global_y_por_clase}, \texttt{Gemini-3.8-Flash} supera a \texttt{GPT-4.1-mini} en las 9 clases, con máximas brechas en \texttt{Further Reading} ($0,9301\text{--}0,9440$ vs. $0,6838\text{--}0,7487$), \texttt{Basis} ($0,5338\text{--}0,5776$ vs. $0,2019\text{--}0,2759$), \texttt{Gap} ($0,4162\text{--}0,4600$ vs. $0,1651\text{--}0,2727$), \texttt{Application} ($0,7690\text{--}0,7753$ vs. $0,5833\text{--}0,6598$) y \texttt{Background} ($0,6368\text{--}0,6951$ vs. $0,4590\text{--}0,5561$). En \textit{0-Shot}, la \textbf{Config B} supera a la Config A en 6 de las 9 clases en ambos modelos. Aunque \textit{9-Shot} en Gemini eleva \texttt{Improvement} ($0,4561 \to 0,5503$) y \texttt{Originator} ($0,4715 \to 0,5243$), reduce las clases mayoritarias (\texttt{Background}: $0,6951 \to 0,6368\text{--}0,6472$; \texttt{Comparison}: $0,6685 \to 0,6195$).

\begin{table}[!t]
\centering
\caption{F1-Score ($\uparrow$) por Clase Canónica en las 8 Variantes de A5 ($N=2.798$)}
\label{tab:f1_por_clase_a5}
\scriptsize
\setlength{\tabcolsep}{1.8pt}
\renewcommand{\arraystretch}{1.04}
\resizebox{\columnwidth}{!}{%
\begin{tabular}{l c c c c c c c c}
\toprule
\textbf{Clase ($n$)} & \textbf{G-0A} $\uparrow$ & \textbf{G-0B} $\uparrow$ & \textbf{G-9A} $\uparrow$ & \textbf{G-9B} $\uparrow$ & \textbf{Gm-0A} $\uparrow$ & \textbf{Gm-0B} $\uparrow$ & \textbf{Gm-9A} $\uparrow$ & \textbf{Gm-9B} $\uparrow$ \\
\midrule
\texttt{Background} ($1.200$) & $0,505$ & $0,556$ & $0,459$ & $0,555$ & $0,692$ & $\mathbf{0,695}$ & $0,647$ & $0,637$ \\
\texttt{Application} ($637$) & $0,649$ & $0,660$ & $0,583$ & $0,584$ & $\mathbf{0,775}$ & $0,775$ & $0,769$ & $0,773$ \\
\texttt{Comparison} ($298$) & $0,635$ & $0,626$ & $0,586$ & $0,568$ & $\mathbf{0,669}$ & $0,657$ & $0,620$ & $0,622$ \\
\texttt{Further Read.} ($227$) & $0,684$ & $0,695$ & $0,744$ & $0,749$ & $0,930$ & $0,941$ & $0,940$ & $\mathbf{0,944}$ \\
\texttt{Originator} ($105$) & $0,477$ & $0,448$ & $0,393$ & $0,340$ & $0,439$ & $0,472$ & $\mathbf{0,524}$ & $0,486$ \\
\texttt{Evidence} ($100$) & $0,288$ & $0,308$ & $0,304$ & $0,328$ & $\mathbf{0,387}$ & $0,378$ & $0,361$ & $0,351$ \\
\texttt{Basis} ($90$) & $0,233$ & $0,276$ & $0,204$ & $0,202$ & $0,566$ & $\mathbf{0,578}$ & $0,534$ & $0,550$ \\
\texttt{Gap} ($73$) & $0,222$ & $0,273$ & $0,168$ & $0,165$ & $0,416$ & $0,446$ & $0,432$ & $\mathbf{0,460}$ \\
\texttt{Improvement} ($68$) & $0,351$ & $0,357$ & $0,291$ & $0,327$ & $0,435$ & $0,456$ & $\mathbf{0,550}$ & $0,526$ \\
\midrule
\textbf{Macro-F1} ($2.798$) & $0,449$ & $0,466$ & $0,415$ & $0,424$ & $0,590$ & $\mathbf{0,600}$ & $0,597$ & $0,594$ \\
\bottomrule
\end{tabular}%
}
\end{table}

\begin{figure*}[!t]
\centering
\includegraphics[height=0.145\textheight,keepaspectratio]{reports/figures/a5/fig02_desempeno_global_y_por_clase.png}
\caption{Desempeño predictivo en A5 sobre el Gold humano ($N = 2.798$): F1 por clase canónica (izquierda) y métricas globales (derecha). Elaboración propia.}
\label{fig:desempeno_global_y_por_clase}
\end{figure*}

La Figura~\ref{fig:confusiones_criticas_y_seccion_citante} revela: (1) \textbf{Perfiles de confusión:} \texttt{GPT 0s-B} sufre sobre-atribución desde \texttt{Background} hacia \texttt{Basis} ($207$, $17,25\pct{}$) y \texttt{Evidence} ($152$, $12,67\pct{}$), \texttt{Application} $\to$ \texttt{Basis} ($134$, $21,04\pct{}$), \texttt{Originator} $\to$ \texttt{Basis} ($36$, $34,29\pct{}$) y \texttt{Further Reading} $\to$ \texttt{Background} ($70$, $30,84\pct{}$), mientras \texttt{Gemini 0s-B} alcanza alta sensibilidad diagonal en \texttt{Further Reading} ($95\pct{}$), \texttt{Basis} ($89\pct{}$), \texttt{Comparison} ($81\pct{}$) y \texttt{Application} ($80\pct{}$), con fugas moderadas de \texttt{Background} a \texttt{Application} ($125$, $10,42\pct{}$) y \texttt{Comparison} ($114$, $9,50\pct{}$); y (2) \textbf{Sección citante:} \texttt{Gemini 0s-B} mantiene alta exactitud en \textit{Conclusion / Future} ($n=42$: $76,19\pct{}$), \textit{Intro / Prelim} ($n=551$: $74,95\pct{}$), \textit{Method / Setup} ($n=362$: $73,20\pct{}$) y \textit{Related Work} ($n=214$: $68,69\pct{}$), mientras en \texttt{GPT 0-Shot} la Config B aporta $\mathbf{+10,28\text{ pp}}$ en \textit{Related Work} ($41,59\pct{} \to 51,87\pct{}$) y $\mathbf{+4,71\text{ pp}}$ en \textit{Intro / Prelim} ($49,55\pct{} \to 54,26\pct{}$).

\begin{figure*}[!t]
\centering
\includegraphics[height=0.145\textheight,keepaspectratio]{reports/figures/a5/fig03_confusiones_criticas_y_seccion_citante.png}
\caption{Matrices de confusión en régimen 0-Shot Config B (\texttt{GPT-4.1-mini} y \texttt{Gemini-3.8-Flash}) y exactitud por macro-sección citante. Elaboración propia.}
\label{fig:confusiones_criticas_y_seccion_citante}
\end{figure*}

\subsubsection{Pruebas Pareadas de Significancia Estadística (A6)}
Evaluamos los $m=12$ contrastes ortogonales con \textit{Paired Bootstrap} bilateral sobre Macro-F1 ($B = 10.000$, semilla $42$), prueba exacta de \textit{McNemar} ($n_{10}$ vs. $n_{01}$) y ajuste de \textbf{Holm-Bonferroni} ($\alpha = 0,05$, Tabla~\ref{tab:significancia_pareada} y Figura~\ref{fig:forest_plot_significancia_pareada}): (1) \textbf{Gemini vs. GPT:} Gemini supera a GPT entre $+0,1331$ y $+0,1828$ Macro-F1 ($p_{\text{boot}} < 10^{-4}$, $p_{\text{McN}} \le 5,32 \times 10^{-50}$, todos $\text{Holm} = \text{SÍ}$, $597\text{--}700$ vs. $150\text{--}210$ aciertos exclusivos); (2) \textbf{Config B vs. A:} en GPT, la Config B mejora significativamente la exactitud McNemar en \textit{0-Shot} ($114/194$, $+2,86\text{ pp}$, $p_{\text{McN}} = 6,75 \times 10^{-6}$, $\text{Holm} = \text{SÍ}$, $+0,0172$ Macro-F1 con $\text{IC}_{95\pct{}} = [+0,0007, +0,0332]$, $p_{\text{boot}} = 0,0368$) y en \textit{9-Shot} ($167/257$, $+3,21\text{ pp}$, $p_{\text{McN}} = 1,54 \times 10^{-5}$, $\text{Holm} = \text{SÍ}$), y en \texttt{Gemini 0-Shot} suma $+0,0096$ Macro-F1 ($0,5900 \to 0,5996$, $p_{\text{boot}} = 0,122$, $136$ discordancias); y (3) \textbf{9-Shot vs. 0-Shot:} degrada significativamente a GPT en Macro-F1 ($-0,0347$ en A, $p_{\text{boot}} = 7,0 \times 10^{-4}$; $-0,0423$ en B, $p_{\text{boot}} = 1,0 \times 10^{-4}$, ambos $\text{Holm} = \text{SÍ}$) y McNemar ($p_{\text{McN}} \le 5,43 \times 10^{-5}$), y en Gemini es neutro en Macro-F1 ($+0,0075$ en A, $p=0,364$; $-0,0053$ en B, $p=0,505$) pero reduce la exactitud McNemar ($182/93$ en B, $p_{\text{McN}} = 1,12 \times 10^{-7}$).

\begin{table*}[!t]
\centering
\caption{Pruebas Pareadas de Significancia Estadística (Paired Bootstrap $B=10.000$ sobre Macro-F1 y McNemar sobre Exactitud con Corrección Holm-Bonferroni $\alpha=0,05, m=12$)}
\label{tab:significancia_pareada}
\scriptsize
\setlength{\tabcolsep}{3.0pt}
\renewcommand{\arraystretch}{1.04}
\resizebox{\textwidth}{!}{%
\begin{tabular}{l l c c c c c c c c c}
\toprule
\textbf{Eje Factorial} & \textbf{Contraste Pareado ($\text{Base} \to \text{Comp.}$)} & \textbf{F1 Base} $\uparrow$ & \textbf{F1 Comp.} $\uparrow$ & \textbf{$\Delta\text{Macro-F1}$} $\uparrow$ & \textbf{$\text{IC}_{95\pct{}}$ Bootstrap} & \textbf{$p$ (Boot)} & \textbf{Holm (Boot)} & \textbf{$n_{10} / n_{01}$} & \textbf{$p$ (McNemar)} & \textbf{Holm (McN)} \\
\midrule
\multirow{4}{*}{\textbf{Config B vs. A}} & \texttt{GPT 0-Shot (B vs A)} & $0,4493$ & $0,4664$ & $\mathbf{+0,0172}$ & $[+0,0007, +0,0332]$ & $3,68\times 10^{-2}$ & NO & $114 / 194$ & $6,75\times 10^{-6}$ & \textbf{SÍ} \\
& \texttt{GPT 9-Shot (B vs A)} & $0,4146$ & $0,4242$ & $+0,0096$ & $[-0,0060, +0,0255]$ & $2,40\times 10^{-1}$ & NO & $167 / 257$ & $1,54\times 10^{-5}$ & \textbf{SÍ} \\
& \texttt{Gemini 0-Shot (B vs A)} & $0,5900$ & $0,5996$ & $+0,0096$ & $[-0,0025, +0,0218]$ & $1,22\times 10^{-1}$ & NO & $62 / 74$ & $3,46\times 10^{-1}$ & NO \\
& \texttt{Gemini 9-Shot (B vs A)} & $0,5974$ & $0,5943$ & $-0,0031$ & $[-0,0148, +0,0081]$ & $5,92\times 10^{-1}$ & NO & $96 / 84$ & $4,12\times 10^{-1}$ & NO \\
\midrule
\multirow{4}{*}{\textbf{9-Shot vs. 0-Shot}} & \texttt{GPT Config A (9s vs 0s)} & $0,4493$ & $0,4146$ & $\mathbf{-0,0347}$ & $[-0,0557, -0,0139]$ & $7,00\times 10^{-4}$ & \textbf{SÍ} & $325 / 219$ & $6,74\times 10^{-6}$ & \textbf{SÍ} \\
& \texttt{GPT Config B (9s vs 0s)} & $0,4664$ & $0,4242$ & $\mathbf{-0,0423}$ & $[-0,0643, -0,0208]$ & $1,00\times 10^{-4}$ & \textbf{SÍ} & $325 / 229$ & $5,43\times 10^{-5}$ & \textbf{SÍ} \\
& \texttt{Gemini Config A (9s vs 0s)} & $0,5900$ & $0,5974$ & $+0,0075$ & $[-0,0086, +0,0239]$ & $3,64\times 10^{-1}$ & NO & $171 / 106$ & $1,20\times 10^{-4}$ & \textbf{SÍ} \\
& \texttt{Gemini Config B (9s vs 0s)} & $0,5996$ & $0,5943$ & $-0,0053$ & $[-0,0202, +0,0104]$ & $5,05\times 10^{-1}$ & NO & $182 / 93$ & $1,12\times 10^{-7}$ & \textbf{SÍ} \\
\midrule
\multirow{4}{*}{\textbf{Gemini vs. GPT}} & \texttt{0-Shot Config A (Gem vs GPT)} & $0,4493$ & $0,5900$ & $\mathbf{+0,1407}$ & $[+0,1156, +0,1658]$ & $< 10^{-4}$ & \textbf{SÍ} & $150 / 647$ & $4,24\times 10^{-69}$ & \textbf{SÍ} \\
& \texttt{0-Shot Config B (Gem vs GPT)} & $0,4664$ & $0,5996$ & $\mathbf{+0,1331}$ & $[+0,1084, +0,1573]$ & $< 10^{-4}$ & \textbf{SÍ} & $168 / 597$ & $5,17\times 10^{-54}$ & \textbf{SÍ} \\
& \texttt{9-Shot Config A (Gem vs GPT)} & $0,4146$ & $0,5974$ & $\mathbf{+0,1828}$ & $[+0,1586, +0,2075]$ & $< 10^{-4}$ & \textbf{SÍ} & $162 / 700$ & $9,89\times 10^{-75}$ & \textbf{SÍ} \\
& \texttt{9-Shot Config B (Gem vs GPT)} & $0,4242$ & $0,5943$ & $\mathbf{+0,1701}$ & $[+0,1455, +0,1952]$ & $< 10^{-4}$ & \textbf{SÍ} & $210 / 646$ & $5,32\times 10^{-50}$ & \textbf{SÍ} \\
\bottomrule
\end{tabular}%
}
\end{table*}

\begin{figure}[!t]
\centering
\includegraphics[width=0.84\columnwidth]{reports/figures/a5/fig04_forest_plot_significancia_pareada.png}
\caption{Gráfico de bosque (\textit{Forest Plot}) de los 12 contrastes ortogonales ($N = 2.798$, $\text{IC}_{95\pct{}}$ Bootstrap $B = 10.000$, Holm-Bonferroni $\alpha = 0,05$). Elaboración propia.}
\label{fig:forest_plot_significancia_pareada}
\end{figure}

\subsubsection{Calibración Probabilística y Regularización del Símplex \texorpdfstring{$\Delta^8$}{Delta-8}}
En $\Delta^8$ (Tabla~\ref{tab:maestra_resultados}), \texttt{Gemini-3.8-Flash} presenta excelente calibración ($\text{ECE} \in [0,1139; 0,1257]$, $\text{Brier} \in [0,4809; 0,5114]$, sobreconfianza $+0,0804\text{ a }+0,1202$) frente a la sobreconfianza de \texttt{GPT-4.1-mini} ($\text{ECE} \in [0,3036; 0,3360]$, $\text{Brier} \in [0,8188; 0,8531]$, sobreconfianza $+0,3032\text{ a }+0,3333$). En \texttt{Gemini 0s-B}, la confianza en aciertos ($0,8290$) supera en $\Delta\text{Conf} = \mathbf{+0,0922}$ a la de errores ($0,7368$, correlación biserial $\rho = \mathbf{+0,4231}$, $p = 7,21 \times 10^{-242}$), frente a $\Delta\text{Conf} = +0,0024$ ($\rho = +0,0479$) en \texttt{GPT 9-Shot}. Mientras en \textit{0-Shot} las 3 modas principales concentran $59,5\pct{}\text{--}64,8\pct{}$ (GPT) y $36,3\pct{}\text{--}37,4\pct{}$ (Gemini), la síntesis por vecindad retórica en \textit{9-Shot} (Tabla~\ref{tab:vecindad_retorica}) reduce en $\mathbf{4\times}$ los ceros exactos en $\Delta^8$ ($36,3\pct{}\text{--}40,9\pct{} \to \mathbf{9,2\pct{}\text{--}10,8\pct{}}$ en GPT; $26,9\pct{}\text{--}28,9\pct{} \to \mathbf{6,5\pct{}\text{--}7,3\pct{}}$ en Gemini), abatiendo la NLL de $5,68\text{--}6,07$ a $\mathbf{2,05\text{--}2,08}$ en GPT y de $1,23\text{--}1,28$ a $\mathbf{1,0687\text{--}1,0913}$ en Gemini.

\subsection{Eficiencia Operativa, Latencia, Tokens, Prompt Caching y Frontera de Pareto (A7)}

\subsubsection{Integridad Operativa del \texorpdfstring{$99,991\pct{}$}{99,991 por ciento} y Diagnóstico Forense de Esquema}
De las $22.384$ inferencias, $\mathbf{22.382}$ cumplieron el esquema JSON ($\mathbf{99,991\pct{}}$; fallo de $0,0089\pct{}$). Los únicos $2$ fallos ocurrieron en \texttt{Gemini-3.8-Flash} Config A sobre instancias Gold \texttt{Further Reading}: $17.318$ en \textit{0-Shot} ($1.450$ entrada, $125$ salida, $0$ pensamiento, $5.477,04\text{ ms}$) y $17.446$ en \textit{9-Shot} ($4.840$ entrada, $77$ salida, $0$ pensamiento, $12.904,01\text{ ms}$), recibiendo $p_k = 1/9$. Con techo de $2.048$ tokens ($125 / 2.048 = 6,1\pct{}$), la longitud de salida de Gemini es bimodal (clúster compacto $67\text{--}76$ tokens y multilínea $114\text{--}124$ tokens): ambos fallos excedieron en exactamente $\mathbf{+1\text{ token}}$ la cota de su clúster ($125 > 124$ y $77 > 76$), con impacto máximo despreciable ($\Delta\text{Macro-F1} \le +0,000298$, $\Delta\text{Accuracy} = +0,000357$).

\subsubsection{Verbosidad JSON, Latencia y Supresión de Thought Tokens}
\texttt{GPT-4.1-mini} emite un JSON en una línea de $\mathbf{66,1}$ tokens ($p_{50} = 915,9\text{--}1.045,6\text{ ms}$, $p_{95} = 1.145,7\text{--}1.829,8\text{ ms}$, $0$ llamadas $>5\text{ s}$ en Config A), mientras \texttt{Gemini-3.8-Flash} serializa un JSON multilínea de $\mathbf{111,7\text{--}119,1}$ tokens ($\sim 1,8\times$ más tokens, $\rho = +0,775$ con latencia, $p_{50} = 2.292,4\text{--}2.481,1\text{ ms}$, $p_{95} = 4.103,2\text{--}7.612,4\text{ ms}$). En el canal de pensamiento de Gemini (\texttt{LOW}): (1) en \textit{0-Shot}, activó pensamiento en el $\mathbf{43,39\pct{}}$ de Config A ($1.214$ llamadas, $206.706$ tokens, media $170,3$) y bajó al $\mathbf{34,95\pct{}}$ en Config B ($978$ llamadas, $170.096$ tokens, $\mathbf{-19,4\pct{}}$ relativo), mientras en \textit{9-Shot} el efecto ancla suprimió el pensamiento en un $\mathbf{99,5\pct{}}$ ($4$ llamadas [$0,14\pct{}$, $674$ tokens] en \texttt{9s-A} y $7$ [$0,25\pct{}$, $2.277$ tokens] en \texttt{9s-B}, bajando el $p_{95}$ en A de $5.269,3$ a $4.103,2\text{ ms}$); (2) en \textit{0-Shot}, el pensamiento gobierna la latencia ($\rho = \mathbf{+0,5567}$, $p < 10^{-140}$), subiendo el $p_{50}$ de $2.154,9$ a $2.795,7\text{ ms}$ en A ($p = 3,10 \times 10^{-154}$) y de $2.158,6$ a $2.832,5\text{ ms}$ en B ($p = 6,17 \times 10^{-139}$) y duplicando el $p_{95}$ ($3,70\text{--}3,89\text{ s} \to 6,14\text{--}6,17\text{ s}$) sin mejorar la exactitud ($67,13\pct{}$ vs. $68,62\pct{}$ en A, $p = 0,414$; $67,79\pct{}$ vs. $68,74\pct{}$ en B, $p = 0,609$; correlación con confianza $r = -0,3031$, $p = 2,83 \times 10^{-119}$); y (3) los $376.802$ tokens de pensamiento en \textit{0-Shot} ($\$3,75/1\text{M}$) costaron $\$1,4130\text{ USD}$.

\subsubsection{Ahorro Financiero por Prompt Caching y Frontera de Pareto}
El costo total de las $22.384$ llamadas bajó de $\mathbf{\$55,1416\text{ USD}}$ brutos a $\mathbf{\$44,4922\text{ USD}}$ netos, ahorrando $\mathbf{\$10,6495\text{ USD}}$ ($\mathbf{-19,31\pct{}}$, Figura~\ref{fig:tokens_cache_eficiencia_y_costos}). En \texttt{GPT-4.1-mini}, mientras \textit{0-Shot} ($\sim 400\text{--}450$ tokens) queda bajo el umbral de $1.024$ tokens ($1,15\pct{}\text{--}1,63\pct{}$ cacheado), en \textit{9-Shot} el caché acertó en el $\mathbf{99,32\pct{}\text{--}99,43\pct{}}$ de las llamadas y reutilizó el $\mathbf{83,06\pct{}}$ (\texttt{9s-A}) y $\mathbf{84,07\pct{}}$ (\texttt{9s-B}) de la entrada ($r = +0,925$), reduciendo el costo de $\$5,6206$ a $\mathbf{\$2,3037\text{ USD}}$ ($\mathbf{-59,01\pct{}}$) en A y de $\$8,2852$ a $\mathbf{\$3,2482\text{ USD}}$ ($\mathbf{-60,80\pct{}}$) en B y recortando el $p_{50}$ en $161,6\text{--}175,3\text{ ms}$ ($p \le 1,27 \times 10^{-3}$, Mann-Whitney). En Gemini \textit{9-Shot}, el caché implícito acertó en $11,29\pct{}$ (\texttt{9s-A}) y $22,80\pct{}$ (\texttt{9s-B}) ($4,33\pct{}$ y $13,39\pct{}$ de tokens), ahorrando $\$2,2647\text{ USD}$ ($3,50\pct{}$ y $11,15\pct{}$). Proyectado a $10.000$ citas (Figura~\ref{fig:frontera_pareto_y_sensibilidad_panel}, izquierda), la \textbf{frontera de Pareto} contiene solo las cuatro variantes \textbf{0-Shot}: (i) \texttt{GPT 0s-A} ($\$5,83/10\text{k}$, $2,62\text{ h}$, $\text{F1}=0,4493$), (ii) \textbf{\texttt{GPT 0s-B}} ($\$6,59/10\text{k}$, $3,11\text{ h}$, $\text{F1}=0,4664$, $\$0,44$ por $+0,01$ F1), (iii) \texttt{Gemini 0s-A} ($\$16,55/10\text{k}$, $7,86\text{ h}$, $\text{F1}=0,5900$), y (iv) el óptimo global \textbf{\texttt{Gemini 0s-B}} ($\$17,68/10\text{k}$, $7,81\text{ h}$, $\text{F1}=\mathbf{0,5996}$, $\$0,79$ por $+0,01$ F1).

\begin{figure*}[!t]
\centering
\includegraphics[height=0.135\textheight,keepaspectratio]{reports/figures/a5/fig08_tokens_cache_eficiencia_y_costos.png}
\caption{Consumo de tokens, costo bruto vs. neto y efecto del caché sobre la latencia $p_{50}$ en régimen 9-Shot ($N = 2.798$ por variante). Elaboración propia.}
\label{fig:tokens_cache_eficiencia_y_costos}
\end{figure*}

\begin{figure*}[!t]
\centering
\includegraphics[height=0.145\textheight,keepaspectratio]{reports/figures/a5/fig11_frontera_pareto_y_sensibilidad_panel.png}
\caption{Frontera de Pareto costo por $10\text{k}$ citas vs. Macro-F1 (izquierda) y exactitud sobre Gold humano por estrato previo A1/A2 (derecha). Elaboración propia.}
\label{fig:frontera_pareto_y_sensibilidad_panel}
\end{figure*}

\section{Discusión y Aportes}

\subsection{Evidencia Local (Config B) frente a Resumen Global (Config A) y Encuadre de Sección}

En respuesta a \textbf{PI1}, nuestros hallazgos extienden empíricamente a \textit{CiteFusion}~\cite{paolini2025citefusion}, \textit{SciRepEval}~\cite{singh2023scirepeval} y Gu et al.~\cite{gu2022local}. Como señalaron Singh et al.~\cite{singh2023scirepeval}, el resumen global (\textbf{Config A}) es estático a nivel documental; cuando una cita ocurre en secciones panorámicas (\textit{Related Work} o \textit{Introduction}), el encabezado citante es genérico y, según el análisis SHAP de \textit{CiteFusion}~\cite{paolini2025citefusion}, induce predicción por exclusión hacia \texttt{Background}. Al inyectar los Top-3 pasajes locales de $\mathcal{P}_{\text{cited}}$ con sus encabezados autorales (\textbf{Config B}, p.~ej., \textit{``6 Fine Estimation''}), el modelo observa directamente si el pasaje citado expone una formulación teórica (\texttt{Basis}), un algoritmo/dataset (\texttt{Application}) o una métrica (\texttt{Comparison}), explicando el salto de $+10,28\text{ pp}$ en \textit{Related Work} en \texttt{GPT-4.1-mini}, el máximo global de $0,5996$ Macro-F1 en \texttt{Gemini-3.8-Flash} y la reducción del $19,4\pct{}$ en pensamiento latente. En respuesta a \textbf{PI2}, las confusiones confirman la triangulación retórica en $\Delta^8$ (\texttt{Basis} $\leftrightarrow$ \texttt{Improvement} $\leftrightarrow$ \texttt{Comparison} $\leftrightarrow$ \texttt{Application}) análoga a \textit{MultiCite}~\cite{lauscher2022multicite}: mientras GPT exhibe \textit{classifier sycophancy}~\cite{fogelson2025llms} sobre-atribuyendo \texttt{Background} hacia \texttt{Basis} ($17,25\pct{}$) y \texttt{Evidence} ($12,67\pct{}$) con $\text{ECE} > 0,30$, Gemini preserva calibración ($\text{ECE} = 0,1158$) y duplica el F1 en \texttt{Basis} ($0,5776$ vs. $0,2759$).

\subsection{Valor Epistemológico de la Curaduría Humana A3, Viabilidad Productiva y Limitaciones}

La Figura~\ref{fig:frontera_pareto_y_sensibilidad_panel} (derecha) y la Tabla~\ref{tab:estratos_a2_vs_humano} demuestran el valor crítico de la curaduría humana (A3) frente a comités 7B--8B (A2). Contra etiquetas pre-A3 ($\mathcal{D}_{\text{test}}^{(\text{panel})}$), el desempeño aparentaba caer en \textit{Needs Review} ($n = 1.120$), sugiriendo un falso ``techo epistémico''. Al evaluar las mismas predicciones contra el Gold humano A3 ($\mathcal{D}_{\text{test}}$), el Macro-F1 de \textbf{\texttt{Gemini 0s-B}} saltó $\mathbf{+0,1510}$ puntos ($0,4485 \to \mathbf{0,5996}$, $+19,98\text{ pp}$ en exactitud y ECE de $0,3156 \to 0,1158$), logrando $\mathbf{69,11\pct{}}$ de exactitud en \textit{Needs Review}, a la par de \textit{Panel Mayoría 2/3} ($70,28\pct{}$) y \textit{Panel Unánime 3/3} ($72,25\pct{}$). El panel 7B--8B acertó la etiqueta humana solo en el $\mathbf{34,24\pct{}}$ ($958 / 2.798$, donde Gemini logra $\mathbf{89,04\pct{}}$); en el $\mathbf{65,76\pct{}}$ donde el panel 7B--8B falló ($n = 1.840$), \textbf{\texttt{Gemini 0s-B} rescata el $\mathbf{57,66\pct{}}$} (vs. $38,64\pct{}$ de \texttt{GPT 0s-B}). Ello explica también por qué \textit{9-Shot} no superó a \textit{0-Shot}: al muestrearse de $\mathcal{D}_{\text{train}}$ (sin curaduría A3), 4 de los 9 ejemplares (\texttt{Background}, \texttt{Improvement}, \texttt{Originator}, \texttt{Further Reading}) provenían del estrato 1-1-1 ($69,73\pct{}$ de error). En respuesta a \textbf{PI3}, \textbf{\texttt{Gemini 0s-B}} ($\$17,68/10\text{k}$) y \textbf{\texttt{GPT 0s-B}} ($\$6,59/10\text{k}$, $p_{95} \le 1,50\text{ s}$) son altamente viables en producción (A8) e inmunes al cambio de prior (\textit{label shift}) entre $\mathcal{D}_{\text{train}}$ ($\texttt{Background} = 12,8\pct{}$) y $\mathcal{D}_{\text{test}}$ ($42,9\pct{}$). Como limitaciones destacan: (i) la exigencia de texto completo en arXiv bajo paridad 1:1; (ii) el techo físico de clases minoritarias (\texttt{Basis}, \texttt{Improvement}); y (iii) la dependencia de umbrales de prefijo ($1.024$ tokens en OpenAI vs. $\ge 4.096$ en Vertex AI) para caché.

\begin{table}[!t]
\centering
\caption{Exactitud ($\uparrow$) sobre el Gold Humano A3 ($N=2.798$) Estratificada por Origen Previo y Concordancia del Panel A2}
\label{tab:estratos_a2_vs_humano}
\scriptsize
\setlength{\tabcolsep}{2.2pt}
\renewcommand{\arraystretch}{1.04}
\resizebox{\columnwidth}{!}{%
\begin{tabular}{l r c c c c}
\toprule
\textbf{Estrato Previo A1/A2 en Test} & \textbf{$n$ ($\pct{}$)} & \textbf{GPT 0s-B} $\uparrow$ & \textbf{GPT 9s-B} $\uparrow$ & \textbf{Gem 0s-B} $\uparrow$ & \textbf{Gem 9s-B} $\uparrow$ \\
\midrule
\textit{Human Gold (MultiCite)} & $563$ ($20,1\pct{}$) & $51,87\pct{}$ & $50,27\pct{}$ & $\mathbf{60,75\pct{}}$ & $56,31\pct{}$ \\
\textit{Panel Unánime (3/3)} & $728$ ($26,0\pct{}$) & $59,89\pct{}$ & $54,81\pct{}$ & $\mathbf{72,25\pct{}}$ & $68,13\pct{}$ \\
\textit{Panel Mayoría (2/3)} & $387$ ($13,8\pct{}$) & $54,78\pct{}$ & $52,97\pct{}$ & $\mathbf{70,28\pct{}}$ & $68,48\pct{}$ \\
\textit{Needs Review (Desacuerdo)} & $1.120$ ($40,0\pct{}$) & $48,66\pct{}$ & $44,82\pct{}$ & $\mathbf{69,11\pct{}}$ & $66,70\pct{}$ \\
\midrule
\textbf{Panel A2 acertó vs. Humano} & $958$ ($34,2\pct{}$) & $80,79\pct{}$ & $73,90\pct{}$ & $\mathbf{89,04\pct{}}$ & $88,20\pct{}$ \\
\textbf{Panel A2 falló vs. Humano} & $1.840$ ($65,8\pct{}$) & $38,64\pct{}$ & $37,01\pct{}$ & $\mathbf{57,66\pct{}}$ & $53,26\pct{}$ \\
\bottomrule
\end{tabular}%
}
\end{table}

\section{Conclusiones y Trabajo Futuro}

Presentamos \textbf{EvidenCite}, un sistema extremo a extremo para recomendación local de evidencia y clasificación retórica de citas bajo las 9 clases de Lyu et al.~\cite{lyu2021classification} en $\Delta^8$. Mediante un pipeline en cinco capas consolidamos $18.658$ instancias con paridad estricta 1:1 entre resumen global (\textbf{Config A}) y Top-3 fragmentos locales recuperados con \textit{SciBERT} y encabezados autorales (\textbf{Config B}), con $\text{leakage}=0$ y un benchmark de prueba de $N = 2.798$ citas curado al $100\pct{}$ por humanos. Sobre $22.384$ inferencias demostramos que \textbf{\texttt{Gemini-3.8-Flash · 0-Shot · Config B}} alcanza el máximo global ($\text{Macro-F1} = 0,5996$, $\text{Weighted-F1} = 0,6933$, $\text{Accuracy} = 68,41\pct{}$, $\text{ECE} = 0,1158$, $\$17,68/10\text{k}$ citas), que la Config B mejora significativamente a \texttt{GPT-4.1-mini} ($+2,86\text{ pp}$ global y $+10,28\text{ pp}$ en \textit{Related Work}), y que la curaduría A3 refuta el falso techo epistémico del panel 7B--8B. Como trabajo futuro se contempla: (1) concluir el entrenamiento en GPU de A4 (\textit{TF-IDF}, \textit{SciBERT} y \textit{PEFT/LoRA}) con calibración de prior (\textit{prior shift}) entre $\mathcal{D}_{\text{train}}$ y $\mathcal{D}_{\text{test}}$ e integrar sus pesos en el servicio cloud de A8 (\texttt{[PENDIENTE\_A8]}); (2) restringir los ejemplares \textit{9-Shot} a citas Gold humanas ($b_{\text{gold}} = \text{True}$) para unir la reducción $4\times$ de ceros en $\Delta^8$ y la supresión del $99,5\pct{}$ de pensamiento con máxima pureza; y (3) escalar el demostrador en AWS ECS Fargate a análisis cienciométricos longitudinales.

\bibliographystyle{IEEEtran}
\bibliography{referencias}

\end{document}
